% Sampled-profile copy: finite-mean shell conditions are retained;
% Section 6 proves the sampled endpoint results directly.
\documentclass[11pt]{article}

\usepackage[letterpaper,margin=1.1in]{geometry}
\usepackage{amsmath,amssymb,amsthm,mathtools}
\usepackage{microtype}
\usepackage{enumitem}
\usepackage[colorlinks=true,allcolors=blue]{hyperref}
\hypersetup{
  pdftitle={Fixed Points and Short Cycles of Luce Permutations: Poisson Limits and Power-Law Endpoint Singularities}
}

\newtheorem{theorem}{Theorem}[section]
\newtheorem{proposition}[theorem]{Proposition}
\newtheorem{lemma}[theorem]{Lemma}
\newtheorem{corollary}[theorem]{Corollary}
\theoremstyle{remark}
\newtheorem{remark}[theorem]{Remark}
\theoremstyle{plain}
\newtheorem{assumption}[theorem]{Assumption}

\newcommand{\E}{\mathbb E}
\newcommand{\Pp}{\mathbb P}
\newcommand{\R}{\mathbb R}
\newcommand{\one}{\mathbf 1}
\newcommand{\cF}{\mathcal F}
\newcommand{\cL}{\mathcal L}
\newcommand{\eps}{\varepsilon}
\newcommand{\Poi}{\operatorname{Poi}}
\newcommand{\Luce}{\operatorname{Luce}}
\newcommand{\TV}{\operatorname{TV}}
\newcommand{\Var}{\operatorname{Var}}
\newcommand{\Fix}{\operatorname{Fix}}

\title{Fixed Points and Short Cycles of Luce Permutations:\\
Poisson Limits and Power-Law Endpoint Singularities}
\author{}
\date{}

\begin{document}
\maketitle

\begin{abstract}
The Luce, or Plackett--Luce, model generates a nonuniform random permutation
by sampling without replacement with probabilities proportional to fixed
weights.  Its global asymptotic geometry was determined in
\cite{borga2025permutonlocallimitsluce}, where a Poisson law for the fixed
points was conjectured.  We prove this conjecture under weak multiscale
profile and endpoint conditions.  The fixed points converge to an inhomogeneous Poisson
process, and the counts of cycles of any fixed collection of lengths converge
jointly to independent Poisson random variables.  Power-law endpoint
singularities produce a second regime in which the cycle means diverge.  We
prove the corresponding joint central limit theorems and locate the cycles on
logarithmic distance scales.  When both endpoints contribute, the cycle
counts arising from them are asymptotically independent and their means add.
\end{abstract}

\tableofcontents
\clearpage

\section{Introduction}

Among nonuniform laws on permutations, the Luce model is one of the simplest
to define.  Give each label a positive weight and repeatedly choose from the
remaining labels with probability proportional to weight.  The rule is the
same at every step and is consistent under deletion of the labels already
chosen.  Equivalently, attach independent exponential clocks to the labels
and order their ringing times.  The first description makes Luce's law a
natural model of sequential choice; the second makes it a tractable random
permutation.  It is therefore a natural setting in which to ask how a
deterministic inhomogeneity changes the familiar microscopic statistics of a
uniform permutation.

The asymptotic study of this object was initiated in
\cite{borga2025permutonlocallimitsluce}.  Starting from a macroscopic weight
profile, that work identified the limiting permuton, computed pattern
densities, proved local limits for consecutive relative patterns, and
established a central limit theorem for inversions.  These results give a
detailed picture of the global geometry and of the relative order seen from a
typical location.  The same work then asked about a statistic that neither
description sees: the number of fixed points.  It predicted a Poisson law
whose mean is the integral of the Luce permuton density along the diagonal
\cite[Section~5.2]{borga2025permutonlocallimitsluce}.  That question is the
starting point of this paper.

The obstacle is not identifying the candidate mean; it is reaching the scale
on which a fixed point lives.  A fixed point is the exact equality of a label
and its rank.  It lies on a lattice-scale diagonal of zero permuton area,
whereas a consecutive-pattern limit records relative order rather than exact
alignment.  Permutations with the same permuton limit can consequently have
completely different fixed-point behavior.  A general cycle criterion in
\cite{mukherjee2016fixed} reaches the correct scale by assuming a uniform
finite-dimensional local law, but the known Luce limits do not imply that
hypothesis.  The gap is most pronounced in the final positions, where few
clocks remain and the limiting density may be singular.

Our first result proves the prediction in
\cite{borga2025permutonlocallimitsluce} under minimal convergence assumptions
on the weight profile and a multiscale condition at the final labels.  We prove
the spatial statement: the fixed points themselves converge to an
inhomogeneous Poisson process, with intensity given by the diagonal of the
Luce permuton density.

A fixed point is a cycle of length one, so the natural next question is
whether the Poisson behavior persists for every fixed cycle length.  It does.
The counts of cycles of lengths $1,\ldots,L$ converge jointly to independent
Poisson variables, with means given by cyclic traces of the permuton density.
This extension is not a formal product of one-point estimates: a cycle is a
closed chain of exact assignments, and its vertices remain coupled through
the common ranking.

Certain structured failures of endpoint tightness reveal a second regime.
Very small weights ring late, while very large weights ring early.  If a
power-law singularity aligns one of these effects with the corresponding
endpoint labels, then every logarithmic distance scale contributes a constant
number of cycles.  The means are then of order
$\log n$.  We prove the corresponding joint central limit theorems, identify
the constants, and locate the cycles on logarithmic distance scales.  The
cycle counts near the left and right endpoints are asymptotically independent,
and their contributions add in the total count.  The size of the cycle
intensity therefore separates two regimes: bounded intensities give Poisson
limits, while diverging intensities give the corresponding centered and
normalized Gaussian limits.  The Sukhatme profile considered in
\cite{borga2025permutonlocallimitsluce} already belongs to the second regime
for fixed points.

The proof changes character at exactly the places suggested by this
description.  In the bulk, concentration of the exponential race identifies
the predictable probability of an exact match.  At the last few positions,
the remaining total weight may be too small to concentrate relatively; a
counting argument avoids that denominator altogether.  For longer cycles we
resample a finite path of clocks into one background permutation, preserving
the combinatorial constraint that each rank has one preimage.  Finally, at a
power-law endpoint we study the surviving population at large times or the
arrived population at small times.  Logarithmic distance turns the resulting
ratio kernel into a translation-invariant one, and the cycle constants become
its return densities.

Let $\theta_{n,1},\ldots,\theta_{n,n}>0$.  A Luce permutation is obtained by
drawing the labels sequentially without replacement, with probability
proportional to their weights.  Thus, if $\pi_n(r)$ denotes the label drawn at
step $r$, then
\begin{equation}\label{eq:luce-law}
 \Pp(\pi_n=\pi)
 =\prod_{r=1}^n
   \frac{\theta_{n,\pi(r)}}
   {\sum_{j=r}^n\theta_{n,\pi(j)}}.
\end{equation}
We use $R_{n,i}=\pi_n^{-1}(i)$ for the rank of label $i$.  The fixed-point
count and process are
\begin{equation}\label{eq:fixed-process}
 X_n:=\sum_{k=1}^n I_{n,k},\qquad
 I_{n,k}:=\one\{R_{n,k}=k\},\qquad
 \Xi_n:=\sum_{k=1}^n I_{n,k}\,\delta_{k/n}.
\end{equation}
This notation separates the draw order $\pi_n$ from the rank map $R_n$; the
two are inverse permutations, while their fixed points agree.

The exponential-race representation is fundamental.  Let
$E_{n,i}\sim\operatorname{Exp}(\theta_{n,i})$ be independent.  Ordering these
variables increasingly gives the Luce draw order, and
\begin{equation}\label{eq:rank-representation}
 R_{n,i}=1+\sum_{j\ne i}\one\{E_{n,j}<E_{n,i}\}.
\end{equation}
For $\ell\ge1$, let
\begin{equation}\label{eq:cycle-counts}
 C_{n,\ell}:=\frac1\ell
 \sum_{(i_1,\ldots,i_\ell)\in[n]_{\ne}^{\ell}}
 \prod_{a=1}^{\ell}\one\{R_{n,i_a}=i_{a+1}\},
 \qquad i_{\ell+1}:=i_1.
\end{equation}
Thus $C_{n,1}=X_n$.  The draw order and rank map are inverses, so they have
the same cycle counts.

\subsection{Context and relation to earlier work}

The model originates in the theory of stochastic choice.  The
proportional-choice axiom was introduced in~\cite{luce1959individual}, and the
induced law on complete rankings was developed in
\cite{plackett1975analysis}.  The construction is closely related to the
Bradley--Terry model for paired comparisons~\cite{bradley1952rank} and to the
independent-utility viewpoint of~\cite{thurstone1927law}.  Probabilistically,
it is successive size-biased sampling without replacement; the exponential
race representation is discussed in~\cite{pitmantran2015sizebiased}.  The
same distribution is the stationary ranking of the Tsetlin library and has
also prompted exact enumerative questions~\cite{chatterjee2024enumerative}.

Cycle counts provide a classical microscopic benchmark.  Under the uniform
law, the counts of cycles of lengths $1,2,\ldots,L$ converge jointly to
independent Poisson variables with means $1,1/2,\ldots,1/L$; the Ewens family
belongs to the same logarithmic-combinatorial framework
\cite{arratia2003logarithmic,tavare2021magical}.  A nonuniform law need not
retain this behavior.  In the fixed-$q$ Mallows model, for example, bounded
cycle counts can have linear means and Gaussian fluctuations, while the
global cycle geometry undergoes a separate localization transition
\cite{gladkich2018cycle,he2023cycles}.  Short cycles are therefore a sensitive
test of which features of the uniform permutation survive inhomogeneity.

Permutons give the natural global limit theory for permutation sequences
\cite{hoppen2013limits}, but their topology averages away exact lattice
coincidences.  A complementary theorem in~\cite{mukherjee2016fixed} shows
that a uniform local approximation to all fixed-dimensional assignment
probabilities implies independent Poisson limits for short cycles, with means
given by cyclic traces of the limiting density.  This theorem explains the
form of our constants, but does not subsume our results.  Its local hypothesis
is stronger than permuton convergence and is not provided by the Luce bulk or
consecutive-pattern limits, especially near a singular endpoint.  We instead
derive microscopic control from the exponential race under the minimal
profile hypotheses below, and then determine what replaces the finite-mean
Poisson law when an endpoint contributes on every logarithmic scale.

\subsection{Finite-mean regime}

Because multiplying all weights by the same constant does not change the
law, for the finite-mean results we normalize
\begin{equation}\label{eq:normalization}
 \frac1n\sum_{i=1}^n\theta_{n,i}=1.
\end{equation}
Let $f_n$ be the step profile
\begin{equation}\label{eq:step-profile}
 f_n(x)=\theta_{n,i},\qquad (i-1)/n<x\le i/n.
\end{equation}

\begin{assumption}\label{ass:profile}
There is a measurable $f:(0,1)\to(0,\infty)$ such that
\begin{equation}\label{eq:L1-profile}
 \lVert f_n-f\rVert_{L^1(0,1)}\longrightarrow0.
\end{equation}
In particular, $\int_0^1f=1$.
\end{assumption}

Equivalently, under the normalization~\eqref{eq:normalization}, it is enough
to assume $f_n\to f$ almost everywhere and $\int f=1$: Scheffe's lemma then
gives~\eqref{eq:L1-profile}.  The integral condition rules out weight mass
escaping through a vanishing fraction of labels.

Define, for $t\ge0$,
\begin{equation}\label{eq:HFD}
 H(t):=\int_0^1e^{-t f(u)}\,du,\qquad
 F(t):=1-H(t),\qquad
 D(t):=\int_0^1 f(u)e^{-t f(u)}\,du.
\end{equation}
Since $f>0$ almost everywhere, $F$ is continuous and strictly increasing from
$0$ to $1$.  Write $t_x:=F^{-1}(x)$ for $0<x<1$.  The limiting Luce permuton
has density
\begin{equation}\label{eq:rho}
 \rho(x,y)=
 \frac{f(x)e^{-f(x)t_y}}{D(t_y)},
 \qquad (x,y)\in(0,1)^2.
\end{equation}

For the finite-mean theorems, the endpoint condition is most naturally
stated on logarithmic depth shells.  Put
\begin{equation}\label{eq:terminal-depth-notation}
 k_{n,m}:=n-m+1,
 \qquad 1\le m\le n.
\end{equation}
For every integer $j\ge1$, let
\begin{equation}\label{eq:terminal-shells}
 \mathcal D_{n,j}:=
 \left\{m:j\le\log\frac{n}{m}<j+1\right\}.
\end{equation}
All shell cutoffs $J$ below tend to infinity through integers.
When $\mathcal D_{n,j}$ is nonempty, define
\begin{equation}\label{eq:shell-floor}
 r_{n,j}:=\max\mathcal D_{n,j},
 \qquad
 b_{n,j}:=\min_{m\in\mathcal D_{n,j}}\theta_{n,k_{n,m}}.
\end{equation}

\begin{assumption}[Endpoint shell capacity]
\label{ass:fixed-endpoint}
Under the normalization~\eqref{eq:normalization},
\begin{equation}\label{eq:shell-endpoint-condition}
 \lim_{J\to\infty}\limsup_{n\to\infty}
 \sum_{\substack{j\ge J\\\mathcal D_{n,j}\ne\varnothing}}
 \exp\{-b_{n,j}j\}=0.
\end{equation}
\end{assumption}

Since $j\le\log(n/r_{n,j})<j+1$, the condition is unchanged if $j$ in
\eqref{eq:shell-endpoint-condition} is replaced by $\log(n/r_{n,j})$ or by
$j+1$.  To see this, split the shells according as $b_{n,j}\le1$ or
$b_{n,j}>1$; the change costs only a fixed factor in the first case and a
summable geometric error in the second.

Condition~\eqref{eq:shell-endpoint-condition} allows the terminal rates to
vanish and permits different lower envelopes on different logarithmic depth
scales.  A convenient stronger sufficient condition is the existence of
$\gamma>0$, $\eps_0>0$, and $n_0$ such that
\begin{equation}\label{eq:endpoint-lower}
 \theta_{n,k}\ge\gamma
 \quad\text{whenever }n\ge n_0\text{ and }(1-\eps_0)n\le k\le n.
\end{equation}
Another useful sufficient condition, which permits the terminal rates to
vanish, is that for some $\delta>0$ and $J_0$, uniformly over $1\le m\le n$
and all sufficiently large $n$,
\begin{equation}\label{eq:logarithmic-endpoint-lower}
 \theta_{n,n-m+1}\ge(1+\delta)
 \frac{\log\log(en/m)}{\log(en/m)}
 \quad\text{whenever }\log(en/m)\ge J_0.
\end{equation}
Indeed, for $m\in\mathcal D_{n,j}$ one has
$\log(en/m)\in[j+1,j+2)$.  Since $x\mapsto(\log x)/x$ is eventually
decreasing, for all sufficiently large $j$,
\[
 b_{n,j}j
 \ge(1+\delta)\frac{j}{j+2}\log(j+2)
 \ge(1+\delta/2)\log(j+2).
\]
Thus the shell costs are bounded by $(j+2)^{-1-\delta/2}$, whose tail is
summable, and \eqref{eq:logarithmic-endpoint-lower} implies
Assumption~\ref{ass:fixed-endpoint}.
The shell condition is strictly weaker and is the hypothesis used in both
finite-mean theorems below.

\begin{theorem}\label{thm:main-poisson}
Under Assumptions~\ref{ass:profile} and~\ref{ass:fixed-endpoint},
\begin{equation}\label{eq:point-process-limit}
 \Xi_n\ \Longrightarrow\ \operatorname{PRM}
 \bigl(\rho(x,x)\,dx\bigr)
\end{equation}
as random finite point measures on $[0,1]$.  In particular,
\begin{equation}\label{eq:count-poisson}
 d_{\TV}\bigl(\cL(X_n),\Poi(\lambda)\bigr)\longrightarrow0,
 \qquad
 \lambda:=\int_0^1\rho(x,x)\,dx<\infty.
\end{equation}
\end{theorem}

The conclusion is spatial, not only numerical: it records where the fixed
points occur.  In particular, it implies the predicted Poisson law for their
total number.

For $\ell\ge1$, define
\begin{equation}\label{eq:cycle-intensity}
 \lambda_\ell:=\frac1\ell\int_{[0,1]^\ell}
 \rho(x_1,x_2)\rho(x_2,x_3)\cdots\rho(x_\ell,x_1)
 \,dx_1\cdots dx_\ell.
\end{equation}

\begin{theorem}\label{thm:short-cycles}
Under Assumptions~\ref{ass:profile} and~\ref{ass:fixed-endpoint}, for every
fixed $L$,
\begin{equation}\label{eq:cycle-vector-limit}
 (C_{n,1},\ldots,C_{n,L}) \Longrightarrow\
 (Z_1,\ldots,Z_L),
\end{equation}
where $Z_1,\ldots,Z_L$ are independent and
$Z_\ell\sim\Poi(\lambda_\ell)$.  The convergence in
\eqref{eq:cycle-vector-limit} also holds in total variation on
$\mathbb Z_+^L$.
\end{theorem}

When $f\equiv1$, one has $\rho\equiv1$ and
$\lambda_\ell=1/\ell$, recovering the classical uniform-permutation law.
The factor $1/\ell$ in~\eqref{eq:cycle-intensity} accounts for the $\ell$
possible choices of a root on a directed cycle.  Both theorems allow a merely
measurable limiting profile and highly irregular triangular arrays; no
interior continuity or uniform upper and lower bounds are imposed.

\subsection{Power-law endpoint regime}

The endpoint condition in the Poisson theorem cannot simply be omitted.  A
power zero at the right endpoint produces cycles near the upper-right corner
of the permutation matrix; a sufficiently strong power pole at the left
endpoint produces them near the lower-left corner.  In either case, every
logarithmic distance scale from the corner contributes a constant amount.
Adding those scales gives order $\log n$ cycles.

The central limit theorems below are the diverging-mean counterparts of the
finite-mean Poisson results.  If $\mu_n\to\mu<\infty$, the corresponding
limit is $\Poi(\mu)$.  If instead $\mu_n\to\infty$, then the natural
comparison is
\[
 \frac{\Poi(\mu_n)-\mu_n}{\sqrt{\mu_n}}
 \Longrightarrow\mathcal N(0,1).
\]
Our endpoint theorems prove the corresponding CLTs for the actual cycle
counts, jointly across cycle lengths, endpoint windows, and active corners.
The relevant mean is of order $\log n$ in the power-law regime and, for fixed
points, $\log\log n$ at the critical pole.

\begin{assumption}\label{ass:simple-power-profile}
For some continuous $f:(0,1)\to(0,\infty)$,
\[
 \theta_{n,i}=f((i-1/2)/n).
\]
At the right endpoint we allow either a power zero
\[
 f(1-s)=c_Rs^\beta(1+O(s^\eta)),
 \qquad \beta,c_R,\eta>0,
\]
or a finite positive limit.  At the left endpoint we allow either a
nonintegrable power pole
\[
 f(s)=c_Ls^{-\alpha}(1+O(s^\eta)),
 \qquad \alpha>1,\quad c_L,\eta>0,
\]
or a finite positive limit.  At least one of the two power behaviors is
present.
\end{assumption}

Call an endpoint \emph{active} when the corresponding power behavior is
present.  Thus $f$ is bounded above and below on every interval separated
from the active endpoints, and there are no other singularities.

Near an active corner, transition probabilities depend asymptotically on the
ratio of two endpoint distances.  Taking logarithms turns this ratio into a
difference.  The resulting increment density is, for $\gamma,q>0$,
\begin{equation}\label{eq:generic-power-kernel}
 p_{\gamma,q}(w):=\gamma q e^{\gamma w}e^{-qe^{\gamma w}},
 \qquad
 \mathfrak b_{\gamma,\ell}(q):=\frac1\ell p_{\gamma,q}^{*\ell}(0).
\end{equation}
Here $p^{*\ell}$ denotes $\ell$-fold convolution.  Thus
$p_{\gamma,q}^{*\ell}(0)/\ell$ is the return density of a closed $\ell$-step
path, divided by the number of choices of its root.
For an active right or left endpoint, respectively, put
\begin{align}
 q_R&:=\Gamma(1+1/\beta)^\beta,
 &b_{R,\ell}&:=\mathfrak b_{\beta,\ell}(q_R),\label{eq:simple-q-right}\\
 q_L&:=\Gamma(1-1/\alpha)^{-\alpha},
 &b_{L,\ell}&:=\mathfrak b_{\alpha,\ell}(q_L).\label{eq:simple-q-left}
\end{align}
Set $b_{\sigma,\ell}=0$ when endpoint $\sigma$ is inactive, and put
$B_\ell:=b_{R,\ell}+b_{L,\ell}$.

The coefficient $b_{\sigma,\ell}$ is the mean number of $\ell$-cycles per
unit of logarithmic distance at corner $\sigma$.  In particular,
$b_{R,1}=\beta q_Re^{-q_R}$ and $b_{L,1}=\alpha q_Le^{-q_L}$.

\begin{theorem}
\label{thm:power-clt}
Under Assumption~\ref{ass:simple-power-profile}, for every fixed $L$,
\begin{equation}\label{eq:power-clt}
 \left(
 \frac{C_{n,\ell}-B_\ell\log n}
 {\sqrt{B_\ell\log n}}
 \right)_{\ell=1}^L
 \ \Longrightarrow\ (Z_1,\ldots,Z_L),
\end{equation}
where $Z_1,\ldots,Z_L$ are independent standard normal variables.  Moreover,
$\E C_{n,\ell}\sim B_\ell\log n$ for every fixed $\ell$.
\end{theorem}

The $\log n$ cycles are not concentrated among the first or last few labels.
They are spread uniformly in logarithmic distance from an active endpoint.
To make this precise, root a cycle $\mathfrak c$ at its largest label
$v(\mathfrak c)$ and put
\[
 d_L(\mathfrak c):=v(\mathfrak c),
 \qquad d_R(\mathfrak c):=n+1-v(\mathfrak c).
\]
Thus $u_\sigma(\mathfrak c):=\log d_\sigma(\mathfrak c)/\log n$ records its
location relative to corner $\sigma$ on a scale from $0$ to $1$.  For
$I=(a,b]\subset(0,1)$, let
\begin{equation}\label{eq:spatial-cycle-count}
 C_{n,\ell}^{\sigma}(I)
 :=\#\{\ell\text{-cycles }\mathfrak c:u_\sigma(\mathfrak c)\in I\}.
\end{equation}
For fixed points this is simply the logarithmic location of the point.

\begin{theorem}
\label{thm:power-spatial}
Under the assumptions of Theorem~\ref{thm:power-clt}, fix pairwise disjoint
intervals $I_1,\ldots,I_J$ whose closures lie in $(0,1)$, and fix $L$.
Jointly over active corners $\sigma$, lengths $1\le\ell\le L$, and
$1\le j\le J$,
\begin{equation}\label{eq:power-spatial-clt}
 \frac{C_{n,\ell}^{\sigma}(I_j)
   -b_{\sigma,\ell}|I_j|\log n}
  {\sqrt{b_{\sigma,\ell}|I_j|\log n}}
 \ \Longrightarrow\ \mathcal N(0,1),
\end{equation}
with independent limits.  Moreover, for every such interval $I$,
\begin{equation}\label{eq:power-spatial-mean}
 \E C_{n,\ell}^{\sigma}(I)
 =b_{\sigma,\ell}|I|\log n+o(\log n).
\end{equation}
\end{theorem}

In particular, the expected number of cycles with
$n^a<d_\sigma\le n^b$ is
$b_{\sigma,\ell}(b-a)\log n+o(\log n)$.  Different logarithmic intervals,
cycle lengths, and endpoints fluctuate independently.  When both endpoints
are active, their means add to the constant
$B_\ell=b_{L,\ell}+b_{R,\ell}$ in Theorem~\ref{thm:power-clt}.

This location describes the entire cycle, not only its root.  More precisely,
for every $\delta>0$, the expected number of cycles counted above that contain
a vertex whose logarithmic distance from the same endpoint differs from
$\log d_\sigma(\mathfrak c)$ by more than $\delta\log n$ tends to zero.

At the borderline pole $f(s)\asymp s^{-1}$, only fixed points diverge, and
they do so on the smaller scale $\log\log n$.

\begin{theorem}
\label{thm:critical-pole}
Suppose $\theta_{n,i}=f((i-1/2)/n)$ for a continuous positive $f$ satisfying
\[
 f(s)=c_Ls^{-1}(1+O(s^\eta))
 \quad(s\downarrow0)
\]
for some $c_L,\eta>0$, and suppose $f$ has a finite positive limit at $1$.
Then
\begin{equation}\label{eq:critical-fixed-clt}
 \frac{C_{n,1}-\log\log n}{\sqrt{\log\log n}}
 \Longrightarrow\mathcal N(0,1),
 \qquad
 \E C_{n,1}\sim\log\log n.
\end{equation}
For every fixed $\ell\ge2$, one has $\sup_n\E C_{n,\ell}<\infty$.
\end{theorem}

Theorem~\ref{thm:critical-pole} is the CLT corresponding to a diverging mean,
now with $\mu_n=\log\log n$.  The means of the longer cycles remain bounded
at the critical exponent.

Section~\ref{sec:power-law} proves these endpoint results directly for
the sampled profiles stated above.  Its population estimates are derived
from the endpoint expansions and the bounds supplied by positivity and
continuity away from the active endpoints.  In particular, the finite-mean
normalization and Assumptions~\ref{ass:profile}--\ref{ass:fixed-endpoint}
are not imposed in the power-law or critical-pole regimes.


The profile singled out in
\cite[Section~5.2]{borga2025permutonlocallimitsluce} gives a concrete instance
of the endpoint regime.

% Lean: Luce.Sukhatme.corollary19 (Luce/Sukhatme.lean).
% Exact statement and constants: Luce/SukhatmeContract.lean;
% the displayed decimal approximation is not part of the Lean contract.
\begin{corollary}\label{cor:sukhatme}
Let $\pi_n$ be a Luce permutation with the standard Sukhatme weights
\[
 \theta_{n,i}=n-i+1,
 \qquad 1\le i\le n.
\]
Put $p(w):=e^we^{-e^w}$ and
$b_\ell:=p^{*\ell}(0)/\ell$.  For every fixed $L$,
\begin{equation}\label{eq:sukhatme-cycle-clt}
 \left(
 \frac{C_{n,\ell}-b_\ell\log n}{\sqrt{b_\ell\log n}}
 \right)_{\ell=1}^L
 \Longrightarrow (Z_1,\ldots,Z_L),
\end{equation}
where the coordinates on the right are independent standard normal random
variables, and $\E C_{n,\ell}\sim b_\ell\log n$.  The spatial conclusion of
Theorem~\ref{thm:power-spatial} also holds, with a single active right corner
and coefficients $b_\ell$.  In particular,
\[
 \frac{C_{n,1}-e^{-1}\log n}{\sqrt{e^{-1}\log n}}
 \Longrightarrow\mathcal N(0,1),
 \qquad
 b_2=K_0(2)\approx0.113894,
\]
where $K_0$ is the modified Bessel function of the second kind.
\end{corollary}

\begin{proof}
Take $f(x)=1-x$ and use the interior grid $i/(n+1)$.
Its weights are $(n-i+1)/(n+1)$, so scale invariance gives precisely the
Sukhatme law.  Lemma~\ref{lem:sp-interior-grid} extends the sampled-profile
conclusions to this grid.  Here $\beta=1$ and $q_R=1$, giving the asserted
constants and spatial limit.  In particular $p(0)=e^{-1}$, and
\[
 \frac12p^{*2}(0)
 =\frac12\int_{\R}e^{-e^w-e^{-w}}\,dw=K_0(2).\qedhere
\]
\end{proof}

Thus the standard Sukhatme model is already in the diverging-mean regime for
fixed points, with mean asymptotic to $e^{-1}\log n$.
Corollary~\ref{cor:sukhatme} proves the corresponding fixed-point CLT and
confirms the behavior suggested by the simulations in
\cite[Section~5.2]{borga2025permutonlocallimitsluce}.  More generally, when a
right power zero and a left nonintegrable power pole are both present, cycles
occur near both endpoints and their contributions add.

\section{Predictable fixed-point probabilities}

We expose the permutation in draw order.  At step $k$, only one event can
create a fixed point: label $k$ must still be available and must be chosen
next.  Its conditional probability is explicit even though the fixed-point
indicators themselves are dependent.

Let $\cF_{n,k}=\sigma(\pi_n(1),\ldots,\pi_n(k))$.  Immediately before draw
$k$, define the remaining weight
\begin{equation}\label{eq:remaining-weight}
 W_{n,k}:=\sum_{i=1}^n\theta_{n,i}\one\{R_{n,i}\ge k\}.
\end{equation}
The conditional chance of a fixed point is
\begin{equation}\label{eq:predictable-p}
 p_{n,k}:=\Pp(I_{n,k}=1\mid\cF_{n,k-1})
 =\frac{\theta_{n,k}\one\{R_{n,k}\ge k\}}{W_{n,k}}.
\end{equation}
Both the numerator indicator and the denominator are
$\cF_{n,k-1}$-measurable.  The random measure
\begin{equation}\label{eq:compensator-measure}
 A_n:=\sum_{k=1}^n p_{n,k}\delta_{k/n}
\end{equation}
is the predictable compensator of $\Xi_n$ in draw order.  The fixed-point
problem is therefore reduced to showing that $A_n$ approaches a deterministic
measure and that its largest atom vanishes.  The following criterion records
this reduction.

\begin{lemma}\label{lem:predictable-poisson}
Let $I_{n,k}\in\{0,1\}$ be adapted, let $x_{n,k}$ be deterministic points in
a compact space $\mathsf X$, and let
$p_{n,k}=\E(I_{n,k}\mid\cF_{n,k-1})$.  Suppose that for a finite deterministic
measure $\nu$ on $\mathsf X$,
\begin{equation}\label{eq:poisson-criterion}
 \sum_kp_{n,k}\delta_{x_{n,k}}\Longrightarrow\nu
 \quad\text{in probability},
 \qquad
 \max_kp_{n,k}\longrightarrow0
 \quad\text{in probability}.
\end{equation}
Then $\sum_k I_{n,k}\delta_{x_{n,k}}$ converges to a Poisson random measure
with intensity $\nu$.
\end{lemma}

\begin{proof}
Fix a nonnegative continuous $g$, abbreviate
$g_{n,k}=g(x_{n,k})$ and $q_{n,k}=1-e^{-g_{n,k}}$.  We first stop the array
before either its total compensator exceeds $K$ or one conditional
probability exceeds $\delta<1$.  Formally, define the predictable indicators
recursively by
\[
 H_{n,k}:=\one\left\{p_{n,k}\le\delta,
 \ \sum_{j<k}H_{n,j}p_{n,j}+p_{n,k}\le K\right\}.
\]
Set
$\widetilde I_{n,k}=H_{n,k}I_{n,k}$ and
$\widetilde p_{n,k}=H_{n,k}p_{n,k}$.  Then
\begin{equation}\label{eq:likelihood-martingale}
 L_{n,m}:=\prod_{k\le m}
 \frac{e^{-g_{n,k}\widetilde I_{n,k}}}
 {1-\widetilde p_{n,k}q_{n,k}}.
\end{equation}
Each factor has conditional mean one, so $(L_{n,m})_m$ is a nonnegative
mean-one martingale.  This likelihood martingale replaces the product
factorization available for independent indicators: here the future
probabilities may depend on the preceding outcomes.

For a single factor $Y_{n,k}$ in~\eqref{eq:likelihood-martingale},
\[
 \E(Y_{n,k}^2\mid\cF_{n,k-1})
 =1+\frac{\widetilde p_{n,k}(1-\widetilde p_{n,k})q_{n,k}^2}
 {(1-\widetilde p_{n,k}q_{n,k})^2}
 \le \exp(C_{g,\delta}\widetilde p_{n,k}).
\]
Iteration therefore gives
\begin{equation}\label{eq:likelihood-L2}
 \sup_m\E L_{n,m}^2
 \le \exp(C_{g,\delta}K),
\end{equation}
so the stopped likelihoods are uniformly integrable.  Meanwhile,
\begin{align*}
 \prod_k(1-\widetilde p_{n,k}q_{n,k})
 &=\exp\left\{-\sum_k\widetilde p_{n,k}q_{n,k}
       +O\left(\max_k\widetilde p_{n,k}
                    \sum_k\widetilde p_{n,k}\right)\right\}\\
 &\longrightarrow
 \exp\left\{-\int(1-e^{-g})\,d\nu\right\}
\end{align*}
in probability whenever no term is deleted.  By~\eqref{eq:poisson-criterion},
choose $K>\nu(\mathsf X)+1$ and any fixed $\delta\in(0,1)$.  The probability
that a term is deleted tends to zero.  Multiplying the last convergence by
the uniformly integrable likelihood in~\eqref{eq:likelihood-martingale}, and
then removing the stopping, gives
\[
 \E\exp\left\{-\sum_k g(x_{n,k})I_{n,k}\right\}
 \longrightarrow
 \exp\left\{-\int(1-e^{-g})\,d\nu\right\},
\]
the Laplace functional of $\operatorname{PRM}(\nu)$.
\end{proof}

\section{The interior compensator}

On every compact subinterval of $[0,1)$, a law of large numbers for the
exponential race controls both the time of the $k$th arrival and the total
rate still present at that time.  This identifies the compensator in the
bulk.  The endpoint, where relative concentration can fail, is postponed to
Section~\ref{sec:endpoint}.

Let
\begin{align}
 \widehat F_n(t)&:=\frac1n\sum_{i=1}^n\one\{E_{n,i}\le t\},
 &F_n(t)&:=\E\widehat F_n(t)
 =\frac1n\sum_{i=1}^n(1-e^{-\theta_{n,i}t}),\label{eq:empirical-F}\\
 \widehat D_n(t)&:=\frac1n\sum_{i=1}^n
 \theta_{n,i}\one\{E_{n,i}\ge t\},
 &D_n(t)&:=\E\widehat D_n(t)
 =\frac1n\sum_{i=1}^n\theta_{n,i}e^{-\theta_{n,i}t}.
 \label{eq:empirical-D}
\end{align}
Let $\tau_{n,k}$ be the $k$th order statistic of the exponential variables.
Then $W_{n,k}=n\widehat D_n(\tau_{n,k})$ and
\begin{equation}\label{eq:survival-indicator}
 \one\{R_{n,k}\ge k\}=\one\{E_{n,k}\ge\tau_{n,k}\}.
\end{equation}

\begin{lemma}\label{lem:race-law}
Under Assumption~\ref{ass:profile}, for every finite $T$,
\begin{equation}\label{eq:uniform-race}
 \sup_{0\le t\le T}|\widehat F_n(t)-F(t)|\longrightarrow0,
 \qquad
 \sup_{0\le t\le T}|\widehat D_n(t)-D(t)|\longrightarrow0
\end{equation}
in probability.  Consequently, for every $\alpha<1$,
\begin{equation}\label{eq:uniform-quantile}
 \max_{1\le k\le\alpha n}|\tau_{n,k}-t_{k/n}|\longrightarrow0
 \quad\text{in probability}
\end{equation}
and
\begin{equation}\label{eq:uniform-denominator}
 \max_{1\le k\le\alpha n}
 \left|\frac{W_{n,k}}n-D(t_{k/n})\right|\longrightarrow0
 \quad\text{in probability}.
\end{equation}
\end{lemma}

\begin{proof}
First, $L^1$ convergence implies
\begin{equation}\label{eq:max-weight}
 \frac1n\max_i\theta_{n,i}\longrightarrow0.
\end{equation}
Indeed, the mass of any interval of length $1/n$ under $f$ tends uniformly to
zero, while $\lVert f_n-f\rVert_1\to0$.

The maps $a\mapsto e^{-ta}$ and $a\mapsto ae^{-ta}$ are uniformly Lipschitz
in $a\ge0$ when $t$ ranges over a compact interval; for the second map one may
use $|(1-ta)e^{-ta}|\le1$.  Hence $F_n\to F$ and $D_n\to D$ uniformly on
$[0,T]$.  At a fixed time, Hoeffding's inequality gives concentration of
$\widehat F_n$, while
\[
 \Var(\widehat D_n(t))
 \le \frac1{n^2}\sum_i\theta_{n,i}^2e^{-\theta_{n,i}t}
 \le \frac{\max_i\theta_{n,i}}n.
\]
A finite grid, monotonicity in $t$, and uniform continuity of $F$ and $D$
upgrade the fixed-time bounds to~\eqref{eq:uniform-race}.

Since $F$ is strictly increasing, uniform convergence of the distribution
functions implies uniform convergence of their quantiles on $[0,\alpha]$,
which gives~\eqref{eq:uniform-quantile}.  Substituting the random quantiles in
the second part of~\eqref{eq:uniform-race} gives
\eqref{eq:uniform-denominator}.
\end{proof}

\begin{proposition}\label{prop:interior-compensator}
Under Assumption~\ref{ass:profile}, for every $\alpha<1$ and every continuous
$g:[0,\alpha]\to\R$,
\begin{equation}\label{eq:weighted-compensator}
 \sum_{k\le\alpha n}g(k/n)p_{n,k}
 \longrightarrow
 \int_0^\alpha g(x)\rho(x,x)\,dx
\end{equation}
in probability.  Moreover,
\begin{equation}\label{eq:interior-max-p}
 \max_{k\le\alpha n}p_{n,k}\longrightarrow0
 \quad\text{in probability}.
\end{equation}
\end{proposition}

\begin{proof}
Put $t_k=t_{k/n}$ and $B_{n,k}=\one\{E_{n,k}\ge t_k\}$.  Since
$d_\alpha:=\min_{0\le t\le t_\alpha}D(t)>0$, Lemma~\ref{lem:race-law} gives
\begin{equation}\label{eq:p-first-approx}
 \sum_{k\le\alpha n}g(k/n)p_{n,k}
 =\frac1n\sum_{k\le\alpha n}
 \frac{g(k/n)\theta_{n,k}\one\{E_{n,k}\ge\tau_{n,k}\}}
 {D(t_k)}+o_{\Pp}(1).
\end{equation}
We next replace the survival indicator by $B_{n,k}$.

Fix $\eta>0$.  The contribution of $k\le\eta n$ is bounded by a constant
times $n^{-1}\sum_{k\le\eta n}\theta_{n,k}$, whose limit is
$\int_0^\eta f$.  For $\eta n<k\le\alpha n$, $t_k$ is bounded away from zero.
On the event in~\eqref{eq:uniform-quantile}, the two indicators can differ
only when $E_{n,k}$ lies in a shrinking interval around $t_k$.  Uniformly on
this range,
\[
 \theta_{n,k}\Pp(|E_{n,k}-t_k|\le\delta)
 \le\int_{t_k-\delta}^{t_k+\delta}
       \theta_{n,k}^2e^{-\theta_{n,k}s}\,ds
 \le C_\eta\delta.
\]
First send $n\to\infty$, then $\delta\downarrow0$, and finally
$\eta\downarrow0$.  This proves that the replacement costs $o_{\Pp}(1)$.

The variables $B_{n,k}$ are independent.  On $k>\eta n$, the variance of the
resulting normalized sum is $O_\eta(n^{-1})$, because
$\sup_{a\ge0}a^2e^{-a t_\eta}<\infty$.  The initial $\eta n$ terms are again
controlled by $n^{-1}\sum_{k\le\eta n}\theta_{n,k}$.  Therefore the sum is
asymptotic in probability to its expectation.  It remains to identify that
expectation.  On $[\eta,\alpha]$, view the sum as the integral of a step
function and put $k_n(x)=\lceil nx\rceil$.  Uniformly there,
$t_{k_n(x)/n}\to t_x$.  The map $a\mapsto ae^{-at}$ is globally
one-Lipschitz, while, for $t\ge t_\eta/2$,
$\sup_{a\ge0}a^2e^{-at}<\infty$.  Hence the uniform convergence of the times,
the continuity of $g$ and $D$, and $\lVert f_n-f\rVert_1\to0$ give
\begin{align*}
 \frac1n\sum_{k\le\alpha n}
 \frac{g(k/n)\theta_{n,k}e^{-\theta_{n,k}t_k}}{D(t_k)}
 &\longrightarrow
 \int_0^\alpha
 \frac{g(x)f(x)e^{-f(x)t_x}}{D(t_x)}\,dx\\
 &=\int_0^\alpha g(x)\rho(x,x)\,dx.
\end{align*}
The discarded initial block is bounded by
$C_\alpha\int_0^\eta f_n(x)\,dx$; it vanishes after first taking
$n\to\infty$ and then $\eta\downarrow0$.  This also handles $x=0$.

For the maximum, on the high-probability event from
Lemma~\ref{lem:race-law},
\[
 \max_{k\le\alpha n}p_{n,k}
 \le \frac{2\max_i\theta_{n,i}}{nd_\alpha}=o(1)
\]
by~\eqref{eq:max-weight}.
\end{proof}

\begin{corollary}\label{cor:interior-poisson}
For every $\alpha<1$,
\begin{equation}\label{eq:interior-poisson}
 \Xi_n|_{[0,\alpha]}\Longrightarrow
 \operatorname{PRM}\bigl(\rho(x,x)\one_{[0,\alpha]}(x)\,dx\bigr).
\end{equation}
\end{corollary}

\begin{proof}
Apply Lemma~\ref{lem:predictable-poisson} and
Proposition~\ref{prop:interior-compensator}.
\end{proof}

\section{The last few positions}\label{sec:endpoint}

It remains to rule out fixed points in the last $o(n)$ positions, uniformly
enough to let $\alpha\uparrow1$.  Directly concentrating the remaining total
weight is fragile: near the end its expectation may be tiny, so an absolute
concentration estimate gives no useful relative error.  We instead condition
on the ringing time of the candidate label and count how many other clocks
survive.  The argument has no random denominator.

For this section only, suppress the first subscript $n$.  Let
\begin{equation}\label{eq:mean-survivors}
 S(t):=\sum_{i=1}^ne^{-\theta_i t}=\E\sum_{i=1}^n\one\{E_i>t\}.
\end{equation}

\begin{lemma}\label{lem:rank-integral}
For every label $k$,
\begin{equation}\label{eq:rank-integral}
 \Pp(R_k=k)=
 \int_0^\infty \theta_k e^{-\theta_k t}
 \Pp\left(\sum_{j\ne k}\one\{E_j>t\}=n-k\right)dt.
\end{equation}
\end{lemma}

\begin{proof}
Condition on $E_k=t$.  The rank is $k$ exactly when $n-k$ of the other
exponential variables are larger than $t$.
\end{proof}

For $m=1,\ldots,n$, write $k_m=n-m+1$ and put
\[
 N(t):=\sum_{j=1}^n\one\{E_j>t\},
 \qquad
 N_{-k}(t):=\sum_{j\ne k}\one\{E_j>t\}.
\]
For $x>0$, also write
\begin{equation}\label{eq:Q-function}
 \mathcal Q(x):=-\log(1-e^{-x}).
\end{equation}

The following block estimate is the basic endpoint-capacity inequality.  It
uses the actual survivor curve $S$, and therefore can be substantially
stronger than a bound obtained only from the mean-one normalization.

\begin{proposition}[Block endpoint capacity]
\label{prop:block-endpoint-capacity}
Let $\mathcal B\subseteq[n]$ be a nonempty set of terminal depths, and put
\[
 r_{\mathcal B}:=\max\mathcal B,
 \qquad
 a_{\mathcal B}:=\min_{m\in\mathcal B}\theta_{k_m}.
\]
For $t_{\mathcal B}>0$, set
\[
 B_{\mathcal B}:=S(t_{\mathcal B})-1.
\]
If $B_{\mathcal B}>r_{\mathcal B}$, then
\begin{equation}\label{eq:block-endpoint-capacity}
 \E\sum_{m\in\mathcal B}I_{k_m}
 \le
 |\mathcal B|
 \exp\left\{-\frac{(B_{\mathcal B}-r_{\mathcal B})^2}
                         {2B_{\mathcal B}}\right\}
 +\mathcal Q(a_{\mathcal B}t_{\mathcal B}).
\end{equation}
\end{proposition}

\begin{proof}
For $0\le t\le t_{\mathcal B}$ and $m\in\mathcal B$,
\[
 \E N_{-k_m}(t)\ge S(t)-1\ge B_{\mathcal B}.
\]
Since $m-1<r_{\mathcal B}<B_{\mathcal B}$, the Chernoff lower-tail bound
gives
\[
 \Pp(N_{-k_m}(t)=m-1)
 \le
 \exp\left\{-\frac{(B_{\mathcal B}-r_{\mathcal B})^2}
                         {2B_{\mathcal B}}\right\}.
\]
Lemma~\ref{lem:rank-integral}, followed by integration of the exponential
density, bounds the total contribution from $t\le t_{\mathcal B}$ by the
first term in~\eqref{eq:block-endpoint-capacity}.

For the remaining times, independence gives
\begin{align*}
 \Pp(N(t)=m-1)
 &\ge \Pp(E_{k_m}\le t,\,N_{-k_m}(t)=m-1)\\
 &= (1-e^{-\theta_{k_m}t})
    \Pp(N_{-k_m}(t)=m-1).
\end{align*}
Consequently,
\[
 \theta_{k_m}e^{-\theta_{k_m}t}
 \Pp(N_{-k_m}(t)=m-1)
 \le
 \frac{\theta_{k_m}}{e^{\theta_{k_m}t}-1}
 \Pp(N(t)=m-1).
\]
For fixed $t>0$, the map $a\mapsto a/(e^{at}-1)$ is decreasing.  Moreover,
the integers $m-1$, $m\in\mathcal B$, are distinct.  Hence the sum of the
last display over $m\in\mathcal B$ is at most
\[
 \frac{a_{\mathcal B}}{e^{a_{\mathcal B}t}-1}
 \sum_{m\in\mathcal B}\Pp(N(t)=m-1)
 \le \frac{a_{\mathcal B}}{e^{a_{\mathcal B}t}-1}.
\]
Integrating from $t_{\mathcal B}$ to infinity gives
\[
 \int_{t_{\mathcal B}}^\infty
 \frac{a_{\mathcal B}}{e^{a_{\mathcal B}t}-1}\,dt
 =-\log(1-e^{-a_{\mathcal B}t_{\mathcal B}}),
\]
which proves the claim.
\end{proof}

There is a complementary individual estimate for a label whose very small
rate is badly aligned with its prescribed terminal depth.  Let
$A_{m-1}^{(-k_m)}$ denote the sum of the $m-1$ smallest rates among
$\{\theta_i:i\ne k_m\}$, with $A_0^{(-k_1)}=0$.

\begin{lemma}[Individual slow-label charge]
\label{lem:slow-label-charge}
% Lean: Luce/EndpointSlowLabel.lean, individual_slow_label_charge.
For every $m$,
\begin{equation}\label{eq:slow-label-charge}
 \Pp(R_{k_m}=k_m)
 \le
 \frac{\theta_{k_m}}
      {\theta_{k_m}+A_{m-1}^{(-k_m)}}.
\end{equation}
\end{lemma}

\begin{proof}
Immediately before the draw at which $m$ clocks remain, condition on the
surviving set.  If it contains $k_m$, the conditional chance that $k_m$ rings
next is its rate divided by the total surviving rate.  The other $m-1$
survivors have total rate at least $A_{m-1}^{(-k_m)}$.  Since a fixed point at
$k_m$ requires both survival to this stage and selection at the next draw,
taking expectations proves~\eqref{eq:slow-label-charge}.
\end{proof}

For a common block $m=1,\ldots,M$, the key deterministic observation is
\begin{equation}\label{eq:two-candidate}
 \sum_{m=1}^M\one\{N_{-k_m}(t)=m-1\}\le2.
\end{equation}
Indeed, if $N(t)=\sum_j\one\{E_j>t\}=s$, then the displayed event can hold only
for $m=s$ when $E_{k_m}>t$, or for $m=s+1$ when $E_{k_m}\le t$.  Thus, at any
fixed time, at most two terminal labels can have the correct survivor count.

\begin{proposition}\label{prop:uniform-endpoint-bound}
Assume~\eqref{eq:normalization}.  Let $1\le M\le\eps_0n$, and suppose
$\theta_{k_m}\ge\gamma$ for $1\le m\le M$.  Put
\begin{equation}\label{eq:B-definition}
 B:=\sqrt{nM},
\end{equation}
and let $s$ be determined by $S(s)=B$.  If $M\le(B-1)/2$ and $s\ge1/\gamma$,
then for a universal $c>0$,
\begin{equation}\label{eq:tail-explicit}
 \E\sum_{m=1}^M I_{k_m}
 \le M e^{-cB}+2e^{-\gamma s}
 \le M e^{-cB}+2\left(\frac{2B}{n}\right)^{\gamma/2}.
\end{equation}
Consequently,
\begin{equation}\label{eq:tail-epsilon}
 \limsup_{n\to\infty}
 \E\sum_{k>(1-\eps)n}I_{n,k}
 \le 2^{1+\gamma/2}\eps^{\gamma/4}
\end{equation}
for all sufficiently small $\eps>0$.
\end{proposition}

\begin{proof}
For $t\le s$ and every $m\le M$,
\[
 \E N_{-k_m}(t)\ge S(t)-1\ge B-1\ge2M.
\]
A Chernoff lower-tail bound for a sum of independent Bernoulli variables gives
\[
 \Pp(N_{-k_m}(t)=m-1)\le e^{-cB}.
\]
Using Lemma~\ref{lem:rank-integral} and integrating the exponential density,
the total contribution of $t\le s$ is at most $Me^{-cB}$.

For $t\ge s\ge1/\gamma$, the map $a\mapsto ae^{-at}$ is decreasing on
$[\gamma,\infty)$.  Hence, by~\eqref{eq:two-candidate},
\begin{align*}
 &\sum_{m=1}^M\int_s^\infty
 \theta_{k_m}e^{-\theta_{k_m}t}
 \Pp(N_{-k_m}(t)=m-1)\,dt\\
 &\hspace{1in}\le
 \int_s^\infty\gamma e^{-\gamma t}
 \sum_{m=1}^M\Pp(N_{-k_m}(t)=m-1)\,dt
 \le2e^{-\gamma s}.
\end{align*}

It remains to bound $s$.  The normalization implies that at least $n/2$
weights are at most $2$.  Thus
\[
 S(t)\ge\frac n2e^{-2t}.
\]
Since $S(s)=B$, we get $s\ge\frac12\log(n/(2B))$, proving the second
inequality in~\eqref{eq:tail-explicit}.

Finally take $M=\lceil\eps n\rceil$.  Then $B/n\to\sqrt\eps$, the first
term in~\eqref{eq:tail-explicit} vanishes, and the required side conditions
hold when $\eps$ is sufficiently small.  This gives~\eqref{eq:tail-epsilon}.
\end{proof}

The particular survivor buffer used in the preceding argument is convenient
but not intrinsic.  A square-root buffer will be useful for both fixed points
and short cycles.

\begin{lemma}[Diagonal shell buffer]
\label{lem:diagonal-shell-buffer}
Under Assumption~\ref{ass:fixed-endpoint}, the choice $h_j=\sqrt j$ satisfies
\begin{equation}\label{eq:buffered-shell-condition}
 \lim_{J\to\infty}\limsup_{n\to\infty}
 \sum_{\substack{j\ge J\\\mathcal D_{n,j}\ne\varnothing}}
 \mathcal Q\bigl(b_{n,j}(j-h_j)\bigr)=0.
\end{equation}
Moreover, $h_j\uparrow\infty$, $h_j<j/4$ eventually, and $j-h_j$ is
nondecreasing.
\end{lemma}

\begin{proof}
Write $x_{n,j}=b_{n,j}j$ and omit empty shells.  Put
\[
 y_{n,j}:=b_{n,j}(j-h_j)
 =x_{n,j}\left(1-\frac1{\sqrt j}\right).
 \]
If $x_{n,j}\le\sqrt j$, then $y_{n,j}\ge x_{n,j}-1$ and hence
$e^{-y_{n,j}}\le e e^{-x_{n,j}}$.  Otherwise,
$y_{n,j}>\sqrt j-1$.  Assumption~\ref{ass:fixed-endpoint} therefore gives
\[
 \lim_{J\to\infty}\limsup_n\sum_{j\ge J}e^{-y_{n,j}}
 \le
 e\lim_{J\to\infty}\limsup_n\sum_{j\ge J}e^{-x_{n,j}}
 +\lim_{J\to\infty}\sum_{j\ge J}e^{1-\sqrt j}=0.
\]
The same estimate shows that $y_{n,j}$ tends uniformly to infinity in the
iterated tail.  Hence $\mathcal Q(y_{n,j})\le2e^{-y_{n,j}}$ there, proving
\eqref{eq:buffered-shell-condition}.  The remaining assertions are immediate
from $h_j=\sqrt j$.
\end{proof}

\begin{proposition}[Shell endpoint tightness]
\label{prop:endpoint-bound}
Under the normalization~\eqref{eq:normalization} and
Assumption~\ref{ass:fixed-endpoint},
\begin{equation}\label{eq:shell-tail-tightness}
 \lim_{J\to\infty}\limsup_{n\to\infty}
 \E\sum_{m\le ne^{-J}} I_{n,k_{n,m}}=0.
\end{equation}
Equivalently,
\begin{equation}\label{eq:shell-tail-epsilon}
 \lim_{\eps\downarrow0}\limsup_{n\to\infty}
 \E\sum_{k>(1-\eps)n}I_{n,k}=0.
\end{equation}
\end{proposition}

\begin{proof}
Let $h_j$ be supplied by Lemma~\ref{lem:diagonal-shell-buffer}.  For a
nonempty shell $\mathcal D_{n,j}$, apply
Proposition~\ref{prop:block-endpoint-capacity} with
 \[
 \mathcal B=\mathcal D_{n,j},
 \qquad
 t_{\mathcal B}=\log\frac{n}{r_{n,j}}-h_j.
\]
For all sufficiently large $j$ this time is positive.  Jensen's inequality
and~\eqref{eq:normalization} give
\[
 S(t_{\mathcal B})
 \ge ne^{-t_{\mathcal B}}
 =r_{n,j}e^{h_j}.
\]
It follows that
$B_{\mathcal B}\ge r_{n,j}e^{h_j}-1>r_{n,j}$ for all sufficiently large
$j$, so the hypothesis of Proposition~\ref{prop:block-endpoint-capacity}
holds.  Moreover, for a universal $c>0$,
\[
 \frac{(B_{\mathcal B}-r_{n,j})^2}{2B_{\mathcal B}}
 \ge c r_{n,j}e^{h_j}.
\]
Thus the first term in
\eqref{eq:block-endpoint-capacity} is at most
\[
 |\mathcal D_{n,j}|e^{-c r_{n,j}e^{h_j}}
 \le r_{n,j}e^{-c r_{n,j}e^{h_j}}.
\]
The nonempty shells have distinct maxima.  Since $h_j$ is nondecreasing,
\[
 \sum_{j\ge J}r_{n,j}e^{-c r_{n,j}e^{h_j}}
 \le\sum_{r\ge1}r e^{-c r e^{h_J}}
 \longrightarrow0
 \qquad(J\to\infty),
\]
uniformly in $n$.  The sum of the second terms in
\eqref{eq:block-endpoint-capacity} tends to zero by
\eqref{eq:buffered-shell-condition}, because
$t_{\mathcal B}\ge j-h_j$ and $\mathcal Q$ is decreasing.  Since the union
of the shells $j\ge J$
is precisely the set of depths with $m\le ne^{-J}$, this proves
\eqref{eq:shell-tail-tightness}; the equivalent formulation
\eqref{eq:shell-tail-epsilon} follows by monotonicity.
\end{proof}

\begin{corollary}[Exceptional-shell refinement]
\label{cor:exceptional-shell-refinement}
% Lean: Luce/EndpointExceptionalTheorem.lean, corollary47_endpoint and corollary47_general.
Assume the normalization~\eqref{eq:normalization}.  For each $n$, let
$\mathcal E_n\subseteq[n]$ be a set of exceptional terminal depths and put
\[
 \mathcal G_{n,j}:=\mathcal D_{n,j}\setminus\mathcal E_n.
\]
When $\mathcal G_{n,j}$ is nonempty, define
\[
 r^{\mathrm G}_{n,j}:=\max\mathcal G_{n,j},
 \qquad
 b^{\mathrm G}_{n,j}
 :=\min_{m\in\mathcal G_{n,j}}\theta_{n,k_{n,m}}.
\]
Let $A_{n,m-1}^{(-k_{n,m})}$ be the sum of the $m-1$ smallest rates after
deleting $\theta_{n,k_{n,m}}$.  Suppose
\begin{align}
 \lim_{J\to\infty}\limsup_{n\to\infty}
 \bigg[
 \sum_{\substack{j\ge J\\\mathcal G_{n,j}\ne\varnothing}}
   e^{-b^{\mathrm G}_{n,j}j}
 +
 \sum_{\substack{m\in\mathcal E_n\\m\le ne^{-J}}}
 \frac{\theta_{n,k_{n,m}}}
      {\theta_{n,k_{n,m}}+A_{n,m-1}^{(-k_{n,m})}}
 \bigg]=0.
 \label{eq:exceptional-shell-condition}
\end{align}
Then the endpoint conclusions
\eqref{eq:shell-tail-tightness}--\eqref{eq:shell-tail-epsilon} hold.
Consequently, Theorem~\ref{thm:main-poisson} remains valid if
Assumption~\ref{ass:fixed-endpoint} is replaced by
\eqref{eq:exceptional-shell-condition}.
\end{corollary}

\begin{proof}
Apply the proof of Lemma~\ref{lem:diagonal-shell-buffer} to the regular
blocks $\mathcal G_{n,j}$.  In the proof of
Proposition~\ref{prop:endpoint-bound}, replace $\mathcal D_{n,j}$,
$r_{n,j}$, and $b_{n,j}$ by their regular-block versions and use the cutoff
$\log(n/r^{\mathrm G}_{n,j})-h_j$.  This cutoff is at least $j-h_j$, so the
same decreasing-$\mathcal Q$ bound makes the regular contribution vanish.
Lemma~\ref{lem:slow-label-charge} bounds the exceptional contribution by the
second sum in~\eqref{eq:exceptional-shell-condition}.  Adding the two proves
the endpoint bounds.  These bounds replace each invocation of
Proposition~\ref{prop:endpoint-bound} in the proof of
Theorem~\ref{thm:main-poisson}, which is otherwise unchanged.
\end{proof}

\begin{remark}\label{rmk:minimal-tail}
The uniform terminal support condition~\eqref{eq:endpoint-lower} implies
Assumption~\ref{ass:fixed-endpoint}.  Indeed, for all sufficiently large
$j$, every label in $\mathcal D_{n,j}$ lies in that terminal interval, so
$b_{n,j}\ge\gamma$ and the shell sum is bounded by the geometric tail
$\sum_{j\ge J}e^{-\gamma j}$.  Proposition
\ref{prop:uniform-endpoint-bound} records the corresponding direct estimate.

At the shell level, \eqref{eq:logarithmic-endpoint-lower} can be weakened
further.  For example,
it is enough that, uniformly on all sufficiently late nonempty shells,
\begin{equation}\label{eq:bertrand-shell-envelope}
 b_{n,j}j
 \ge \log j+(1+\delta)\log\log j
\end{equation}
for some $\delta>0$, because the shell costs are then bounded by
$1/(j(\log j)^{1+\delta})$.  In terms of an individual terminal depth with
$q=\log(n/m)$, the corresponding envelope has order
\[
 \theta_{n,n-m+1}
 \ge
 \frac{\log q+(1+\delta)\log\log q}{q}.
\]
The usual iterated-log refinements of the convergent Bertrand series give
still weaker sufficient envelopes, and sparse bad shells need only have
summable costs.  Provided the remaining shell costs already have uniformly
vanishing tails, a single moving anomalous shell is allowed whenever its
product $b_{n,j}j$ tends to infinity.

The minimum $b_{n,j}$ can nevertheless be spoiled by one harmless ultra-slow
label.  Corollary~\ref{cor:exceptional-shell-refinement} removes such labels
from their regular shells and charges them individually.  This exception
mechanism is proved here only for fixed points.  The short-cycle argument
below uses the lower floor on every label in a shell, so no analogue under
the exceptional-shell refinement is claimed.
\end{remark}

\begin{proof}[Proof of Theorem~\ref{thm:main-poisson}]
For fixed $\alpha<1$, Corollary~\ref{cor:interior-poisson} gives the desired
Poisson limit on $[0,\alpha]$.  Proposition~\ref{prop:endpoint-bound} and
Markov's inequality imply
\begin{equation}\label{eq:tail-tightness}
 \lim_{\alpha\uparrow1}\limsup_{n\to\infty}
 \Pp\bigl(\Xi_n((\alpha,1])>0\bigr)=0.
\end{equation}
It remains to verify that the proposed limiting intensity is finite; weak
convergence alone does not imply convergence of means.  Fix
$\alpha<\beta<1$.  Interior count convergence, followed by truncation and
Fatou's lemma, gives
\begin{align*}
 \int_\alpha^\beta\rho(x,x)\,dx
 &=\E\Poi\left(\int_\alpha^\beta\rho(x,x)\,dx\right)\\
 &\le\liminf_{n\to\infty}\E\Xi_n((\alpha,\beta])
 \le\limsup_{n\to\infty}\E\Xi_n((\alpha,1]).
\end{align*}
Letting $\beta\uparrow1$ and using monotone convergence gives
\[
 \int_\alpha^1\rho(x,x)\,dx
 \le\limsup_{n\to\infty}\E\Xi_n((\alpha,1]).
\]
Proposition~\ref{prop:endpoint-bound} shows that the right side tends to zero
as $\alpha\uparrow1$, first along $\alpha=1-e^{-J}$ and then for arbitrary
$\alpha$ by monotonicity.  Together with finiteness of the intensity on every
compact subinterval of $[0,1)$, this proves both
\begin{equation}\label{eq:intensity-tail-bound}
 \lambda:=\int_0^1\rho(x,x)\,dx<\infty,
 \qquad
 \int_\alpha^1\rho(x,x)\,dx\longrightarrow0.
\end{equation}
Now let $\alpha\uparrow1$ in the interior point-process limit.  Equations
\eqref{eq:tail-tightness} and~\eqref{eq:intensity-tail-bound} give the full
convergence~\eqref{eq:point-process-limit}, and hence
$X_n\Rightarrow\Poi(\lambda)$.

For integer-valued random variables, convergence in distribution to a proper
integer-valued law implies pointwise convergence of the probability mass
functions; Scheffe's lemma then upgrades this to total variation convergence,
giving~\eqref{eq:count-poisson}.
\end{proof}

\section{Short cycles}\label{sec:short-cycles}

This section proves Theorem~\ref{thm:short-cycles}.  There are again two
tasks.  In the bulk, a linear reservoir of moderate-rate clocks yields the
finite-dimensional local law needed for factorial moments.  Near the
endpoint, the fixed-point estimate no longer suffices: the label that closes
a late cycle may come from anywhere.  We handle this by inserting every
proposed finite path into a common background permutation.  The construction
retains the essential permutation constraint that each rank has one
preimage.

For an integer $m$, write $[m]^r_{\ne}$ for the set of ordered $r$-tuples of
distinct elements of $[m]$.

\begin{lemma}\label{lem:moderate-reservoir}
Under Assumption~\ref{ass:profile}, for every $\alpha<1$ there are
$d,\eta>0$ such that, for all sufficiently large $n$,
\begin{equation}\label{eq:moderate-reservoir}
 \#\{i:\theta_{n,i}\ge d\}\ge(\alpha+2\eta)n.
\end{equation}
Consequently, after deleting a fixed number of clocks and before $\alpha n$
arrivals, the remaining total rate is at least $d\eta n$.
\end{lemma}

\begin{proof}
Since $f>0$ almost everywhere, choose $d,\eta>0$ so that
$|\{x:f(x)\ge2d\}|>\alpha+3\eta$.  The part of this set on which $f_n<d$ has
measure at most $d^{-1}\lVert f_n-f\rVert_1$, which proves
\eqref{eq:moderate-reservoir}.  At most $\alpha n+O(1)$ members of the
reservoir have arrived or been deleted, leaving at least $\eta n$ clocks of
rate at least $d$.
\end{proof}

\begin{lemma}\label{lem:cyclic-local}
Under Assumption~\ref{ass:profile}, fix $r\ge1$, $\alpha<1$, a permutation
$\tau\in S_r$, and a continuous function $g:[0,\alpha]^r\to\R$.  Then
\begin{align}
 &\sum_{\boldsymbol i\in[\lfloor\alpha n\rfloor]^r_{\ne}}
 g(\boldsymbol i/n)
 \Pp\bigl(R_{n,i_a}=i_{\tau(a)},\ 1\le a\le r\bigr)
 \label{eq:cyclic-local-limit}\\
 &\hspace{0.7in}\longrightarrow
 \int_{[0,\alpha]^r}g(\boldsymbol x)
 \prod_{a=1}^r\rho(x_a,x_{\tau(a)})\,d\boldsymbol x.
 \notag
\end{align}
Moreover, uniformly over distinct marked labels $i_1,\ldots,i_r$ and
distinct ranks $j_1,\ldots,j_r\le\alpha n$,
\begin{equation}\label{eq:weighted-bulk-cylinder}
 \Pp(R_{n,i_a}=j_a,\ 1\le a\le r)
 \le K_{r,\alpha}n^{-r}\prod_{a=1}^r\theta_{n,i_a}.
\end{equation}
\end{lemma}

\begin{proof}
Delete the marked clocks and denote the unmarked order statistics by
$T_0=0<T_1<\cdots<T_{n-r}$.  If $W_q^\circ$ is the remaining unmarked rate,
the memoryless property gives
\begin{equation}\label{eq:race-gap-representation}
 T_{q+1}-T_q=\frac{\xi_q}{W_q^\circ},
\end{equation}
where, conditionally on the elimination chain, the $\xi_q$ are independent
standard exponentials.  Lemma~\ref{lem:moderate-reservoir} gives
$W_q^\circ\ge d\eta n$ for $q\le\alpha n$.  Sorting the prescribed ranks,
and grouping marked clocks which must enter the same unmarked gap, therefore
gives
\[
 \Pp(R_{n,i_a}=j_a,\ a\le r)
 \le K_{r,\alpha}n^{-r}\prod_{a=1}^r\theta_{n,i_a},
\]
which is \eqref{eq:weighted-bulk-cylinder}.

We identify the asymptotic first with all marked rates bounded by a fixed
$M$.  Deleting finitely many such clocks changes the empirical race by
$o(1)$, uniformly in their labels.  Lemma~\ref{lem:race-law} and
\eqref{eq:race-gap-representation} then show, uniformly when the prescribed
ranks are separated by at least $\zeta n$, that
\begin{equation}\label{eq:bounded-marked-asymptotic}
 n^r\Pp(R_{n,i_a}=j_a,\ a\le r)
 -\prod_{a=1}^r
 \frac{\theta_{n,i_a}e^{-\theta_{n,i_a}t_{j_a/n}}}
 {D(t_{j_a/n})}
 \longrightarrow0.
\end{equation}
Indeed, each marked clock is inserted into a relevant unmarked gap; Taylor
expansion of its exponential density is uniform for rates at most $M$, and
the rescaled separated gaps converge jointly to independent standard
exponentials.

The contribution to \eqref{eq:cyclic-local-limit} of tuples containing a
rate larger than $M$ is, by \eqref{eq:weighted-bulk-cylinder}, at most a
constant times
\begin{equation}\label{eq:bulk-high-truncation}
 \sum_{a=1}^r
 \left(\frac1n\sum_i\theta_{n,i}\one\{\theta_{n,i}>M\}\right)
 \left(\frac1n\sum_i\theta_{n,i}\right)^{r-1}.
\end{equation}
It vanishes as first $n\to\infty$ and then $M\to\infty$.  After this rate
truncation, tuples with two coordinates within $\zeta n$ have contribution
$O_{r,M}(\zeta)$; the discarded high-rate part is still controlled by
\eqref{eq:bulk-high-truncation}.

It remains to pass from the grid kernel in
\eqref{eq:bounded-marked-asymptotic} to $\rho$.  On rates at most $M$, the
kernel is bounded and uniformly continuous in its rank coordinate on
$[0,\alpha]$, so the grid sum differs from the corresponding step integral by
$o_M(1)$.  The map $a\mapsto ae^{-at}$ is one-Lipschitz, uniformly for
$a,t\ge0$, and $D(t_y)$ is bounded below for $y\le\alpha$.  A telescoping
product and $\lVert f_n-f\rVert_1\to0$ identify the integral.  The parts on
which either $f_n$ or $f$ exceeds $M$ are bounded by their $L^1$ tails,
because every numerator is at most its source rate.  Let $n\to\infty$, then
$\zeta\downarrow0$, and finally $M\to\infty$ to obtain
\eqref{eq:cyclic-local-limit}.
\end{proof}

We now set up the endpoint resampling argument.  Fix a proposed $\ell$-cycle
and abbreviate $\theta_{n,i}$ to $\theta_i$.  Generate a background family
$E_i^0\sim\operatorname{Exp}(\theta_i)$, with order statistics
\[
 0=T_0<T_1<\cdots<T_n<T_{n+1}=\infty,
\]
and replace the clocks on the proposed cycle by independent copies.  Put
\begin{equation}\label{eq:ghost-window}
 J_j:=\bigl(T_{(j-\ell)\vee0},T_{(j+\ell)\wedge(n+1)}\bigr),
 \qquad
 p_{i,j}:=\int_{J_j}\theta_i e^{-\theta_i t}\,dt.
\end{equation}
The window $J_j$ contains every time that could acquire rank $j$ after these
$\ell$ replacements.  Since replacing $\ell$ clocks changes every rank by at
most $\ell$, averaging over the background gives
\begin{equation}\label{eq:ghost-cylinder-bound}
 \Pp(R_{i_a}=j_a,\ 1\le a\le\ell)
 \le\E\prod_{a=1}^\ell p_{i_a,j_a}.
\end{equation}
Each time belongs to at most $2\ell+1$ windows, and consequently
\begin{equation}\label{eq:ghost-row-bound}
 \sum_{j=1}^np_{i,j}\le2\ell+1.
\end{equation}

The next lemma iterates this insertion bound along a finite path.  Its
distinct-source condition is exactly the condition supplied by a genuine
cycle.

\begin{lemma}\label{lem:finite-insertion-path}
For fixed $\ell,m\ge1$ and every $v\in[n]$,
\begin{equation}\label{eq:finite-insertion-path}
 \E\sum_{\substack{u_0,\ldots,u_{m-1}\in[n]\\
                    u_0,\ldots,u_{m-1}\ {\mathrm{distinct}}}}
 p_{u_0,u_1}\cdots p_{u_{m-2},u_{m-1}}p_{u_{m-1},v}
 \le K_{\ell,m}.
\end{equation}
\end{lemma}

\begin{proof}
Generate one independent replacement $E_i^1$ for every background clock.
For a fixed source sequence $\boldsymbol u$, the product in
\eqref{eq:finite-insertion-path} is the conditional probability that every
$E_{u_a}^1$ falls in the window indexed by the next vertex, ending with
$J_v$.  Swap $E_{u_a}^0$ with $E_{u_a}^1$ for all $a$.  These swaps concern
distinct labels and preserve the joint law.  If $\widetilde R$ is the rank
permutation of the new background, success before the swap implies
\begin{equation}\label{eq:approximate-permutation-path}
 |\widetilde R_{u_a}-u_{a+1}|\le q_{\ell,m}\quad(a<m-1),
 \qquad
 |\widetilde R_{u_{m-1}}-v|\le q_{\ell,m},
\end{equation}
where $q_{\ell,m}=\ell+m+2$ is sufficient: membership in an old window costs
at most $\ell+1$ ranks, and the $m$ swaps cost at most another $m$.

For a deterministic permutation $\sigma$, give $u$ an edge to every $w$ with
$|\sigma(u)-w|\le q_{\ell,m}$.  Each target has indegree at most
$2q_{\ell,m}+1$, since every rank has one preimage.  Thus at most
$(2q_{\ell,m}+1)^m$ length-$m$ paths end at $v$.  Sum the
measure-preserving comparison over $\boldsymbol u$.
\end{proof}

Put $c_j=\sum_kp_{k,j}$.  Expanding this sum and separating whether its source
is new gives the useful consequence
\begin{equation}\label{eq:added-predecessor}
 \sup_{v\in[n]}\E
 \sum_{\substack{i_2,\ldots,i_\ell\in[n]\setminus\{v\}\\
                  i_2,\ldots,i_\ell\ {\mathrm{distinct}}}}
 c_{i_2}p_{i_2,i_3}\cdots p_{i_\ell,v}\le K_\ell.
\end{equation}
Indeed, when the source $k$ in $c_{i_2}=\sum_kp_{k,i_2}$ is distinct from
$i_2,\ldots,i_\ell$, Lemma~\ref{lem:finite-insertion-path} applies.  If
$k=i_a$, discard the extra factor $p_{i_a,i_2}\le1$ and apply the lemma to
the remaining path.  The interpretation for $\ell=1$ is
$\sup_v\E c_v<\infty$.

We first prepare two truncations: high rates globally and low rates on compact
label intervals.  For $M>1$ and $\delta>0$, let
\[
 H_{n,M}:=\{i:\theta_{n,i}>M\},\qquad
 L_{n,\delta}:=\{i:\theta_{n,i}<\delta\},
\]
and let $V_{n,L}(S)$ be the number of vertices in $S$ which belong to cycles
of length at most $L$.

\begin{proposition}[Rate truncation]\label{prop:cycle-rate-truncation}
Under Assumption~\ref{ass:profile},
\begin{align}
 \lim_{M\to\infty}\limsup_n\E V_{n,L}(H_{n,M})&=0,
 \label{eq:high-cycle-truncation}\\
 \lim_{\delta\downarrow0}\limsup_n
 \E V_{n,L}\bigl(L_{n,\delta}\cap[\lfloor\beta n\rfloor]\bigr)&=0
 \qquad(\beta<1).
 \label{eq:compact-low-cycle-truncation}
\end{align}
\end{proposition}

\begin{proof}
The $L^1$ profile convergence gives
\begin{equation}\label{eq:profile-tail-controls}
 \lim_{M\to\infty}\limsup_n\frac1n
 \sum_{i\in H_{n,M}}\theta_{n,i}=0,
 \qquad
 \lim_{\delta\downarrow0}\limsup_n\frac{|L_{n,\delta}|}{n}=0.
\end{equation}
For the second relation, compare $\{f_n<\delta\}$ with
$\{f<2\delta\}$ and use $f>0$ almost everywhere.

Choose $M$ so large that the labels of rate at most $M$ carry total weight at
least $n/2$.  If $\theta_i>M$, then, with
$G_n(t)=\sum_k\theta_k e^{-\theta_kt}$,
\[
 \theta_i e^{-\theta_i t}\le\frac{2\theta_i}{n}G_n(t).
\]
Indeed, all rates at most $M$ have survival factor at least
$e^{-\theta_i t}$.  Since $c_j=\int_{J_j}G_n(t)dt$, this gives
\begin{equation}\label{eq:high-entry-comparison}
 p_{i,j}\le\frac{2\theta_i}{n}c_j.
\end{equation}
Start a proposed cycle at $i\in H_{n,M}$, apply
\eqref{eq:ghost-cylinder-bound} and then
\eqref{eq:high-entry-comparison} to its first edge.  For $\ell\ge2$,
\begin{align*}
 \Pp(i\text{ belongs to an }\ell\text{-cycle})
 &\le \E\sum_{\substack{i_2,\ldots,i_\ell\ne i\\
                         i_2,\ldots,i_\ell\ {\mathrm{distinct}}}}
 p_{i,i_2}p_{i_2,i_3}\cdots p_{i_\ell,i}\\
 &\le \frac{2\theta_i}{n}\E
 \sum_{\substack{i_2,\ldots,i_\ell\ne i\\
                  i_2,\ldots,i_\ell\ {\mathrm{distinct}}}}
 c_{i_2}p_{i_2,i_3}\cdots p_{i_\ell,i}
 \le K_\ell\frac{\theta_i}{n}
\end{align*}
by \eqref{eq:added-predecessor}.  For fixed points, use
$p_{i,i}\le2\theta_i c_i/n$ and the $\ell=1$ interpretation of that estimate.
Summing over $i$ and $\ell\le L$ gives
\begin{equation}\label{eq:quantitative-high-cycles}
 \E V_{n,L}(H_{n,M})
 \le\frac{K_L}{n}\sum_{i\in H_{n,M}}\theta_{n,i}.
\end{equation}
This proves \eqref{eq:high-cycle-truncation}.  Now fix $\beta<1$.  Lemma
\ref{lem:moderate-reservoir}, applied at a fixed number larger than $\beta$,
and \eqref{eq:race-gap-representation} give
\begin{equation}\label{eq:interior-window-length}
 \sup_{i\le\beta n}\E|J_i|\le K_L/n.
\end{equation}
In a cycle avoiding $H_{n,M}$, the edge closing into $i$ is at most
$M|J_i|$.  Sum all other edges with \eqref{eq:ghost-row-bound} to get
\[
 \Pp\bigl(i\text{ lies in a short cycle avoiding }H_{n,M}\bigr)
 \le K_{L,M}/n.
\]
A cycle meeting both sets contains at most $L$ vertices from
$L_{n,\delta}\cap[\lfloor\beta n\rfloor]$ per high-rate vertex.  Therefore
\[
 \E V_{n,L}\bigl(L_{n,\delta}\cap[\lfloor\beta n\rfloor]\bigr)
 \le L\E V_{n,L}(H_{n,M})
   +K_{L,M}\frac{|L_{n,\delta}|}{n}.
\]
Use \eqref{eq:profile-tail-controls}, then let $M\to\infty$, to prove
\eqref{eq:compact-low-cycle-truncation}.
\end{proof}

For $\alpha<1$, let $C_{n,\ell}^{(\alpha)}$ count cycles all of whose labels
are at most $\alpha n$.

\begin{proposition}[Short-cycle shell tightness]
\label{prop:cycle-endpoint-bound}
Under Assumptions~\ref{ass:profile} and~\ref{ass:fixed-endpoint}, for every
fixed $L$,
\begin{equation}\label{eq:cycle-tail-tightness}
 \lim_{\alpha\uparrow1}\limsup_n
 \E\sum_{\ell=1}^L
 \bigl(C_{n,\ell}-C_{n,\ell}^{(\alpha)}\bigr)=0.
\end{equation}
\end{proposition}

\begin{proof}
The case $\ell=1$ is Proposition~\ref{prop:endpoint-bound}.  Fix
$2\le\ell\le L$ and put $B=2\ell+1$.  We first restrict to cycles avoiding
$H_{n,M}$.  Proposition~\ref{prop:cycle-rate-truncation} shows that the
discarded contribution vanishes as $M\to\infty$.

Fix a large shell cutoff $J_0$ and then $\beta<1$ so close to one that, for
all sufficiently large $n$, every label larger than $\beta n$ belongs to a
shell with index at least $J_0$.  We also discard cycles meeting
$L_{n,\delta}\cap[\lfloor\beta n\rfloor]$.  Their contribution vanishes as
$\delta\downarrow0$ by Proposition~\ref{prop:cycle-rate-truncation}.
Thus every remaining cycle vertex has rate at most $M$, and every remaining
vertex at most $\beta n$ has rate at least $\delta$.

Rotate a proposed cycle uniquely as
\[
 u\longrightarrow v\longrightarrow i_3\longrightarrow\cdots
 \longrightarrow i_\ell\longrightarrow u,
 \qquad v=\max\{u,v,i_3,\ldots,i_\ell\}.
\]
Conditional on the common background in \eqref{eq:ghost-window}, let
\[
 A_{v,u}:=\sum_{i_3,\ldots,i_\ell}
 p_{v,i_3}\cdots p_{i_\ell,u},
\]
where the intermediate sums range over retained labels; both distinctness
and the genuine-cycle restrictions $i_a<v$ are dropped, while the outer pair
remains restricted by $u<v$.  For $\ell=2$ this means
$A_{v,u}=p_{v,u}$.  Iterating
\eqref{eq:ghost-row-bound} gives the pointwise bound
\begin{equation}\label{eq:marked-return-bound}
 \sum_u A_{v,u}\le B^{\ell-1}.
\end{equation}
Consequently, for any set $S$ of retained predecessors and every $t$,
\begin{equation}\label{eq:marked-edge-bound}
 \sum_v\one\{t\in J_v\}\sum_{u\in S}
 \theta_u e^{-\theta_ut}A_{v,u}
 \le B^\ell\sup_{u\in S}\theta_u e^{-\theta_ut},
\end{equation}
because at most $B$ target windows contain $t$.
Here the supremum is zero when $S$ is empty, and enlarging the actual sum
from $u<v$ to all pairs can only increase its left side.
The unique rotation cancels the factor $1/\ell$ in the definition of the
cycle count.  Hence \eqref{eq:ghost-cylinder-bound} gives
\[
 \E\bigl[\#\{\text{retained $\ell$-cycles}\}\bigr]
 \le \E\sum_{u<v}p_{u,v}A_{v,u},
\]
with any further restrictions imposed on the outer pair in both sums.

Suppose first that the terminal depth of $v$ belongs to
$\mathcal D_{n,j}$, $j\ge J$.  The part of $p_{u,v}$ coming from
$t<j-h_j$ can be nonzero only if
\[
 T_{n-r_{n,j}+1-\ell}<j-h_j.
\]
Indeed, $v\ge n-r_{n,j}+1$ and $J_v$ starts at $T_{v-\ell}$.  Jensen's
inequality and the normalization give
\[
 S(j-h_j)\ge ne^{-j+h_j}\ge r_{n,j}e^{h_j}.
\]
For all large $j$, a Chernoff lower-tail bound therefore yields
\[
 \Pp\bigl(T_{n-r_{n,j}+1-\ell}<j-h_j\bigr)
 \le \exp\{-c r_{n,j}e^{h_j}\}.
\]
Pointwise in the common background, the early part of $p_{u,v}$ is at most
the indicator of this event, so the total rooted weight for each $v$ is at
most $B^{\ell-1}$ times that indicator.  Taking expectations and using
$|\mathcal D_{n,j}|\le r_{n,j}$, all early-window contributions are bounded
by a constant times
\begin{equation}\label{eq:early-cycle-shells}
 \sum_{\substack{j\ge J\\\mathcal D_{n,j}\ne\varnothing}}
 r_{n,j}e^{-c r_{n,j}e^{h_j}}
 \le\sum_{r\ge1}r e^{-c r e^{h_J}}
 \longrightarrow0.
\end{equation}
Here the maxima $r_{n,j}$ of the nonempty shells are distinct integers and
$h_j\ge h_J$ for $j\ge J$.

It remains to integrate $p_{u,v}$ over $t\ge j-h_j$.  If $u>\beta n$, its
terminal depth belongs to a shell $\mathcal D_{n,r}$ with $r\ge J_0$.
Since $u<v$, one has $r\le j$, and Lemma
\ref{lem:diagonal-shell-buffer} gives
$t\ge j-h_j\ge r-h_r$.  By increasing $J_0$, we may make the limsup in
\eqref{eq:buffered-shell-condition} smaller than $\mathcal Q(1)$; then, for
all sufficiently large $n$, every nonempty source shell $r\ge J_0$ satisfies
$b_{n,r}(r-h_r)>1$.  As
$b_{n,r}\le\theta_u\le M$, the map $a\mapsto ae^{-at}$ is decreasing on
this range, and hence
\[
 \theta_u e^{-\theta_u t}\le b_{n,r}e^{-b_{n,r}t}.
\]
Apply \eqref{eq:marked-edge-bound} separately to each source shell,
retaining $u<v$ (and hence $r\le j$) before enlarging the time range to
$[r-h_r,\infty)$, and then sum over the shells.  Their contribution is
at most
\begin{equation}\label{eq:late-cycle-shells}
 B^\ell\sum_{\substack{r\ge J_0\\\mathcal D_{n,r}\ne\varnothing}}
 \int_{r-h_r}^\infty b_{n,r}e^{-b_{n,r}t}\,dt
 =B^\ell\sum_{\substack{r\ge J_0\\\mathcal D_{n,r}\ne\varnothing}}
 e^{-b_{n,r}(r-h_r)},
\end{equation}
which vanishes as $J_0\to\infty$ by
\eqref{eq:buffered-shell-condition}.

Finally, if $u\le\beta n$, then $\delta\le\theta_u\le M$.  For $J$ large,
$\delta(J-h_J)>1$, so on $t\ge j-h_j\ge J-h_J$,
\[
 \theta_u e^{-\theta_u t}\le\delta e^{-\delta t}.
\]
The same marked-edge bound makes this contribution at most
\begin{equation}\label{eq:late-cycle-bulk}
 B^\ell\int_{J-h_J}^\infty\delta e^{-\delta t}\,dt
 =B^\ell e^{-\delta(J-h_J)},
\end{equation}
which vanishes as $J\to\infty$.

Choose first $M$, then $J_0$ and the corresponding $\beta$, then $\delta$,
and finally $J$, keeping these cutoffs fixed before taking each limsup in
$n$.  Equations
\eqref{eq:early-cycle-shells}--\eqref{eq:late-cycle-bulk}, followed by the
two truncation estimates, prove \eqref{eq:cycle-tail-tightness}; replacing
$1-e^{-J}$ by an arbitrary $\alpha\uparrow1$ only requires a rounding check.
Integer rounding can add at most one terminal-depth label, which for all
large $n$ lies in a shell with index at least $J-1$.  Thus the estimate with
$J-1$ proves the result along $\alpha=1-e^{-J}$, and monotonicity gives the
general limit.
\end{proof}

\begin{proof}[Proof of Theorem~\ref{thm:short-cycles}]
Fix $\alpha<1$ and nonnegative integers $m_1,\ldots,m_L$, and put
$r=\sum_{\ell=1}^L\ell m_\ell$.  Expanding the joint falling factorial gives
an ordered collection of vertex-disjoint cycles.  Lemma
\ref{lem:cyclic-local} therefore yields
\begin{equation}\label{eq:cycle-factorial-moments}
 \E\prod_{\ell=1}^L(C_{n,\ell}^{(\alpha)})_{m_\ell}
 \longrightarrow
 \prod_{\ell=1}^L\lambda_\ell(\alpha)^{m_\ell},
\end{equation}
where
\[
 \lambda_\ell(\alpha):=\frac1\ell
 \int_{[0,\alpha]^\ell}
 \prod_{a=1}^\ell\rho(x_a,x_{a+1})\,d\boldsymbol x,
 \qquad x_{\ell+1}=x_1.
\]
These are the joint factorial moments of independent Poisson variables, so
the truncated cycle vector has the asserted limit.

Proposition~\ref{prop:cycle-endpoint-bound} gives the full endpoint
tightness~\eqref{eq:cycle-tail-tightness}.

First-moment convergence and \eqref{eq:cycle-tail-tightness} show that
\[
 \sup_{\alpha<\beta<1}
 \bigl[\lambda_\ell(\beta)-\lambda_\ell(\alpha)\bigr]
 \longrightarrow0
 \qquad(\alpha\uparrow1).
\]
Hence the increasing $\lambda_\ell(\alpha)$ have finite
limits, necessarily the integrals in \eqref{eq:cycle-intensity}.  Markov's
inequality removes the label cutoff in \eqref{eq:cycle-factorial-moments} and
proves \eqref{eq:cycle-vector-limit}.  Finally, weak convergence on the
countable lattice $\mathbb Z_+^L$ gives pointwise convergence of mass
functions; Scheffe's lemma upgrades it to total variation convergence.
\end{proof}

\section{Power-law endpoint singularities and cycle CLTs}
\label{sec:power-law}

We prove Theorems~\ref{thm:power-clt} and~\ref{thm:power-spatial}
directly under Assumption~\ref{ass:simple-power-profile}.  The population
estimates needed in the proof are consequences of that assumption, not
additional assumptions on a triangular array.  The critical pole is treated
separately under the hypotheses of Theorem~\ref{thm:critical-pole}.
The mean-one normalization used in the finite-mean part of the paper is not
imposed in this section.

\subsection{Population estimates for sampled profiles}

Write $E_{n,i}\sim\operatorname{Exp}(\theta_{n,i})$ for the independent
clocks and define
\begin{align}
 H_n(t)&=\frac1n\sum_i e^{-t\theta_{n,i}},&
 G_n(t)&=1-H_n(t),&
 D_{n,r}(t)&=\frac1n\sum_i\theta_{n,i}^r e^{-t\theta_{n,i}}.
 \label{eq:sp-transforms}
\end{align}
We abbreviate $D_{n,1}$ to $D_n$.  Fix a sufficiently small
$\epsilon_*>0$.  At each active endpoint, the rates throughout the
$\epsilon_*$-neighborhood are between two fixed positive multiples of the
specified power.  On the complement of the active neighborhoods the rates
are in a fixed interval $[d,M]\subset(0,\infty)$.  These facts follow by
first shrinking the neighborhoods so that their relative errors are less
than $1/2$, and then using continuity, positivity, and the finite positive
limits at inactive endpoints.  No differentiability or monotonicity of $f$
is used.

\begin{lemma}[Population transforms]\label{lem:sp-populations}
Under Assumption~\ref{ass:simple-power-profile}, at an active right endpoint,
put $A_-=\Gamma(1+1/\beta)c_R^{-1/\beta}$.  For some $\zeta>0$,
uniformly as $t\to\infty$ and $nt^{-1/\beta}\to\infty$,
\begin{align}
 H_n(t)&=A_-t^{-1/\beta}
  \left(1+O\left(t^{-\zeta}+(nt^{-1/\beta})^{-1}\right)\right),\notag\\
 D_n(t)&=\frac{A_-}{\beta}t^{-1-1/\beta}
  \left(1+O\left(t^{-\zeta}+(nt^{-1/\beta})^{-1}\right)\right).
 \label{eq:sp-slow-transforms}
\end{align}
At an active left endpoint, put
$A_+=\Gamma(1-1/\alpha)c_L^{1/\alpha}$.  Uniformly as $t\downarrow0$
and $nt^{1/\alpha}\to\infty$,
\begin{align}
 G_n(t)&=A_+t^{1/\alpha}
  \left(1+O\left(t^\zeta+(nt^{1/\alpha})^{-1}\right)\right),\notag\\
 D_n(t)&=\frac{A_+}{\alpha}t^{1/\alpha-1}
  \left(1+O\left(t^\zeta+(nt^{1/\alpha})^{-1}\right)\right).
 \label{eq:sp-fast-transforms}
\end{align}
The constants and exponent may depend on the active endpoint.
After deleting any fixed number $r$ of clocks, the same statements hold
for the sums divided by $n$, with constants uniform in the deleted labels.
\end{lemma}

\begin{proof}
We estimate the unweighted transforms together with $tD_n(t)$.  For the
explicit monotone or unimodal kernels below, midpoint averaging introduces
an error bounded by their supremum times $O(n^{-1})$.

At the right endpoint, reflect the sampling grid and compare $f(1-s)$
with $cs^\beta$, where $c=c_R$.  The pure-power kernels have integrals
\[
 \int_0^\infty e^{-ct s^\beta}\,ds=A_-t^{-1/\beta},\qquad
 \int_0^\infty ct s^\beta e^{-ct s^\beta}\,ds
 =\frac{A_-}{\beta}t^{-1/\beta}.
\]
Their suprema are uniformly bounded, and their tails beyond $s=1$ are
exponentially small for $t\ge1$.

The endpoint expansion and the mean value theorem bound the errors in
both kernels by
\[
 Ct s^{\beta+\eta}e^{-dt s^\beta}
\]
near zero, after absorbing a linear factor into the exponential.  This
unimodal envelope has integral $O(t^{-(1+\eta)/\beta})$ and supremum
$O(t^{-\eta/\beta})$.  Away from the endpoint, both rates are bounded
below, so both kernels are exponentially small; here
$tae^{-ta}\le(2/e)e^{-ta/2}$.  Consequently,
\[
 |H_n(t)-A_-t^{-1/\beta}|
 +\left|tD_n(t)-\frac{A_-}{\beta}t^{-1/\beta}\right|
 \le C\left(t^{-(1+\eta)/\beta}+n^{-1}\right).
\]
Division by the leading terms proves \eqref{eq:sp-slow-transforms}.

At the left endpoint, choose $0<\eta'<\min\{\eta,\alpha-1\}$ and put
$q=\alpha-\eta'>1$.  With $c=c_L$, the pure-power integrals are
\[
 \int_0^\infty(1-e^{-ct s^{-\alpha}})\,ds=A_+t^{1/\alpha},\qquad
 \int_0^\infty ct s^{-\alpha}e^{-ct s^{-\alpha}}\,ds
 =\frac{A_+}{\alpha}t^{1/\alpha}.
\]
Their midpoint errors are $O(n^{-1})$, and their tails beyond $1$ are
$O(t)$.  The perturbation errors near zero are bounded by
\[
 Ct s^{-q}e^{-dt s^{-\alpha}}
 \le Ct s^{-q}e^{-dt s^{-q}},\qquad 0<s<1.
\]
The last kernel has integral $O(t^{1/q})$ and bounded supremum.  Away
from zero, boundedness of both rates gives an $O(t)$ contribution.
Since $t\le t^{1/q}$ for $0<t\le1$,
\[
 |G_n(t)-A_+t^{1/\alpha}|
 +\left|tD_n(t)-\frac{A_+}{\alpha}t^{1/\alpha}\right|
 \le C(t^{1/q}+n^{-1}).
\]
This gives \eqref{eq:sp-fast-transforms} with
$\zeta=1/q-1/\alpha>0$.

Finally, deleting at most $r$ terms changes either unweighted sum divided
by $n$ by at most $r/n$, and changes $tD_n$ by at most $r/(en)$.
These errors are already allowed, uniformly in the deleted labels.
\end{proof}

\begin{lemma}[Quantiles and deterministic weight bounds]
\label{lem:sp-quantiles}
At an active right endpoint define $t^R_{n,m}$ by
$H_n(t^R_{n,m})=m/n$; at an active left endpoint define $t^L_{n,m}$
by $G_n(t^L_{n,m})=m/n$.  There is $\xi>0$ such that, uniformly for
$m\to\infty$ and $m/n\to0$,
\begin{align}
 t^R_{n,m}&=(A_-n/m)^\beta(1+O(e_{n,m})),&
 \theta_{n,n-m+1}t^R_{n,m}&=q_R(1+O(e_{n,m})),\notag\\
 \frac{\theta_{n,n-m+1}}{nD_n(t^R_{n,m})}
  &=\frac{\beta q_R}{m}(1+O(e_{n,m})),
 &&\label{eq:sp-right-quantile}\\
 t^L_{n,m}&=(m/(A_+n))^\alpha(1+O(e_{n,m})),&
 \theta_{n,m}t^L_{n,m}&=q_L(1+O(e_{n,m})),\notag\\
 \frac{\theta_{n,m}}{nD_n(t^L_{n,m})}
  &=\frac{\alpha q_L}{m}(1+O(e_{n,m})),
 &&\label{eq:sp-left-quantile}
\end{align}
where $e_{n,m}=(m/n)^\xi+m^{-1}$ and $q_R,q_L$ are exactly
\eqref{eq:simple-q-right}--\eqref{eq:simple-q-left}.
For $t=t^\sigma_{n,m}$ and $t/4\le s\le4t$,
\begin{equation}
 nD_n(s)\asymp m/t,\qquad nD_{n,2}(s)\le Cm/t^2.
 \label{eq:sp-rough-moments}
\end{equation}
The corresponding survivor mean $nH_n$ on the right, or arrival mean
$nG_n$ on the left, changes by at least $cum$ between $t$ and
$(1\pm u)t$, for $0<u\le1/2$.  Moreover
$D_n((1\pm u)t)/D_n(t)=1+O(u)$.

After deleting at most $r$ clocks, the total rate in a gap with $h$
survivors at the right endpoint is deterministically at least
\begin{equation}
 c n^{-\beta}h^{\beta+1}.
 \label{eq:sp-right-weight-floor}
\end{equation}
At a gap after $h+O(r)$ early arrivals at the left endpoint it is at least
\begin{equation}
 c_r n^\alpha h^{1-\alpha},
 \label{eq:sp-left-weight-floor}
\end{equation}
provided $h\ge4r+4$ and $h\le\epsilon_* n/4$.
For ranks in any fixed compact subinterval of $(0,1)$, the deterministic
quantiles lie in a compact subinterval of $(0,\infty)$, $nD_n$ is
comparable to $n$, and the remaining rate is deterministically at least
$c n$, uniformly after bounded deletions.
\end{lemma}

\begin{proof}
Put $x=m/n$.  Positivity of the rates makes $H_n$ strictly decreasing and
$G_n$ strictly increasing, so the quantiles exist uniquely for $0<m<n$.

We first justify inversion in the joint regime.  On the right, choose a
fixed large $S$ where the power error in Lemma~\ref{lem:sp-populations}
is a small fraction of the leading term.  For $x$ small and $m$ large,
$H_n(S)>x$, hence $t=t^R_{n,m}>S$.  At this quantile the additive
estimate gives
\[
 |x-A_-t^{-1/\beta}|
 \le Ct^{-\eta/\beta}A_-t^{-1/\beta}+C/n.
\]
Absorbing the small errors first gives $t\asymp x^{-\beta}$;
substitution back into the inequality gives the stated relative expansion.
The same argument on the left, starting at a fixed small $S$ with
$G_n(S)>x$, gives $t^L_{n,m}\asymp x^\alpha$ and then its relative
expansion.  Thus the population remainders become positive powers of $x$.

The sampled endpoint expansions have an additional $O(m^{-1})$ error
because the endpoint distance is $x(1-1/(2m))$.  Multiplying by the
quantile expansions gives the rate-time products, using
\[
 c_RA_-^\beta=q_R,\qquad c_LA_+^{-\alpha}=q_L.
\]
The weighted population estimates give
\[
 nt^R_{n,m}D_n(t^R_{n,m})=\frac m\beta(1+O(e_{n,m})),\qquad
 nt^L_{n,m}D_n(t^L_{n,m})=\frac m\alpha(1+O(e_{n,m})).
\]
Choose a common positive $\xi$ below all the resulting error exponents.
Division then proves the formulas for $\theta/(nD_n)$.

For $t=t^\sigma_{n,m}$, Lemma~\ref{lem:sp-populations} applies
uniformly on $[t/8,8t]$ and gives $nD_n(s)\asymp m/t$.  The inequality
\[
 D_{n,2}(s)\le\frac{2}{es}D_n(s/2)
\]
therefore proves the second-moment bound on $[t/4,4t]$.  Integrating
the exact derivatives
\[
 H_n'=-D_n,\qquad G_n'=D_n,\qquad D_n'=-D_{n,2}
\]
gives the mean separation and relative weight stability.  No asymptotic
formula is differentiated.

At the right endpoint, the global bound
\[
 \theta_{n,i}\ge c((n-i+1)/n)^\beta
\]
implies that any $h$ remaining labels have total rate at least
\[
 cn^{-\beta}\sum_{a=1}^h a^\beta\ge c'n^{-\beta}h^{\beta+1}.
\]
At the left endpoint, after at most $2h$ total removals, at least $h$
of the first $3h$ labels remain.  Each has rate at least $c(n/h)^\alpha$,
proving the left floor.  Bounded deletion and displacement budgets are
absorbed into this count by increasing the fixed block multiple if
necessary and shrinking the endpoint neighborhood.

Finally, the population estimates, or global rate bounds at inactive
endpoints, provide fixed $0<a<b$ with $G_n(a)<\delta$ and
$H_n(b)<\delta$.  Thus interior quantiles lie in $[a,b]$.  A fixed
interior block of bounded positive rates gives $D_n\ge c$ there, while
$D_n(s)\le1/(es)$ gives the upper bound.  Bounded deletions preserve
these estimates.  A linear number of survivors has total rate at least
$cn$, by the right floor when that endpoint is active and by a global
positive rate lower bound otherwise.
\end{proof}


\subsection{Joint insertion estimates}

Throughout this subsection we assume Assumption~\ref{ass:simple-power-profile}.
The number of marked labels is fixed; constants may depend on this number
and on a fixed moment order.  No monotonicity of the sampled profile is used.
Write
\[
 \iota_R(m)=n-m+1,\quad \iota_L(m)=m,\qquad
 d_R(i)=n-i+1,\quad d_L(i)=i.
\]
Choose a sufficiently small fixed $\delta>0$.  At an active corner let
\[
 x_R(a,h)=(a/h)^\beta,\qquad x_L(a,h)=(h/a)^\alpha,
 \qquad K_\sigma(a,h)=h^{-1}k_\sigma(a/h),
\]
where
\[
 k_R(u)=\beta q_Ru^\beta e^{-q_Ru^\beta},\qquad
 k_L(u)=\alpha q_Lu^{-\alpha}e^{-q_Lu^{-\alpha}}.
\]
Here $q_R,q_L$ have the values in
\eqref{eq:simple-q-right}--\eqref{eq:simple-q-left}.
For a sufficiently small $c>0$ put
\[
 H_R(a,h)=\frac1h(a/h)^\beta e^{-c(a/h)^\beta},\qquad
 H_L(a,h)=\frac1h(h/a)^\alpha e^{-c(h/a)^\alpha}.
\]
The constant $c$ will always retain exponential slack: decreasing it a
fixed number of times is allowed.  In particular, any fixed polynomial in
$x_\sigma$ can be absorbed by this slack.

\paragraph{The exact marked probability space.}
Delete a set $D$ of at most $r$ labels.  Order the remaining independent
clocks as
\[
 0=T_0^D<T_1^D<\cdots<T_N^D<T_{N+1}^D=\infty,
 \qquad N=n-|D|.
\]
For a label $i\in D$ define its conditional probability of entering gap
$q$ by
\begin{equation}\label{eq:sp-gap-factor}
 X_{i,q}^D=\int_{T_q^D}^{T_{q+1}^D}
                  \theta_{n,i}e^{-\theta_{n,i}t}\,dt.
\end{equation}
Suppose that distinct marked labels $i_1,\ldots,i_s$ are required to have
distinct full ranks $j_1,\ldots,j_s$, and take $D=\{i_1,\ldots,i_s\}$.
If $j_{e_1}<\cdots<j_{e_s}$, set
\begin{equation}\label{eq:sp-demanded-gap}
 q_{e_b}=j_{e_b}-b,\qquad 1\le b\le s.
\end{equation}
For each repeated value $q$, integrate the corresponding marked clocks
over the ordered simplex in $(T_q^D,T_{q+1}^D)$, in the order of their
demanded ranks.  Denote this integral by $Z_q^D$.  Independence of the
marked clocks gives the exact identity
\begin{equation}\label{eq:sp-cylinder-simplex}
 \Pp(R_{n,i_e}=j_e\text{ for all }e)
       =\E\prod_{q}Z_q^D,
 \qquad
 \prod_qZ_q^D\le\prod_eX_{i_e,q_e}^D.
\end{equation}
The second assertion follows by enlarging each ordered simplex to its
product box.  If the $q_e$ are distinct, it is an equality.  This accounts
for adjacent demanded ranks without assigning independent ranks to clocks
which must occupy one ordered gap.

\begin{lemma}[Moments of one insertion]\label{lem:sp-insertion-moments}
Fix positive integers $r,p_0$.  There are
$h_0,\nu_0,\nu,c,C>0$ such that the following assertions hold, uniformly
after deletion of at most $r$ labels and uniformly for gap indices within
$r+1$ of the full-rank gap under consideration.  At an active corner,
for $h_0\le a,h\le\delta n$ and $1\le p\le p_0$, put
\[
 \mathcal M_\sigma(a,h)
       =\{x_\sigma(a,h)\le(a\wedge h)^{\nu_0}\},
 \quad
 E_R(a,h)=e^{-ca^\nu}\one_{\mathcal M_R(a,h)^c},\quad
 E_L(a,h)=e^{-ch^\nu}\one_{\mathcal M_L(a,h)^c}.
\]
Then
\begin{equation}\label{eq:sp-insertion-lp}
 \|X_{\iota_\sigma(a),q}^D\|_p
                  \le C\{H_\sigma(a,h)+E_\sigma(a,h)\}.
\end{equation}
The exceptional kernels have bounded row and column sums.  More precisely,
after decreasing $\nu$ and $c$,
\begin{equation}\label{eq:sp-extreme-restricted-sums}
 \sup_{a\ge A}\sum_{h\ge A}E_\sigma(a,h)
 +\sup_{h\ge A}\sum_{a\ge A}E_\sigma(a,h)
 \le Ce^{-cA^\nu}.
\end{equation}
The same restricted-sum assertion holds for
$H_\sigma\one_{\mathcal M_\sigma^c}$.
\end{lemma}

\begin{proof}
Fix $0<\nu_0\le1$, increase $h_0$ to absorb the deletion and gap-shift
budgets, and shrink $\delta$ as necessary.  Write
\[
 i=\iota_\sigma(a),\qquad \theta=\theta_{n,i},\qquad
 s=T_q^D,\qquad x=x_\sigma(a,h).
\]
Conditional on the background elimination chain, the nonfinal gap length
is $\Delta=\xi/W$, where $\xi$ is standard exponential and independent
of the preceding spacings, hence of $s$.  Lemma~\ref{lem:sp-quantiles}
and endpoint comparability give
\[
 W\ge cn^{-\beta}h^{\beta+1}\ (R),\qquad
 W\ge cn^\alpha h^{1-\alpha}\ (L),\qquad \theta/W\le Cx/h.
\]
Since $X_{i,q}^D=e^{-\theta s}(1-e^{-\theta\Delta})$, integration
over $\xi$ yields
\begin{equation}\label{eq:sp-insertion-moment-reduction}
 \E[(X_{i,q}^D)^p]
 \le p!(Cx/h)^p\,\E e^{-p\theta s}.
\end{equation}

Put $\tau_h^R=(n/h)^\beta$ and $\tau_h^L=(h/n)^\alpha$.  At a
sufficiently small fixed multiple of $\tau_h^\sigma$, the population
estimates place the survivor mean on the right well above the gap
threshold, or the arrival mean on the left well below it.  These
comparisons retain a fixed margin after the allowed deletions and shifts.
Chernoff bounds therefore give
\begin{equation}\label{eq:sp-order-lower-tail}
 \Pp(s<k\tau_h^\sigma)\le Ce^{-dh},\qquad
 \E e^{-p\theta s}\le C(e^{-dx}+e^{-dh}),
\end{equation}
using $\theta\tau_h^\sigma\asymp x$.  On the moderate set, $x\le h$,
so \eqref{eq:sp-insertion-moment-reduction}, followed by taking $p$-th
roots, gives $C(x/h)e^{-cx}$.

For the right extreme region, put $\gamma=\beta/(\beta+1)$.  If
$h\le a^\gamma$, repeat the survivor-count argument at comparison depth
$y=a^\gamma$, that is, at time $k(n/y)^\beta$.  The mean is a
sufficiently large multiple of $y$, and
$\theta(n/y)^\beta\asymp(a/y)^\beta=y$.  Thus the survival expectation
is at most $Ce^{-da^\gamma}$.  If $h>a^\gamma$, use
\eqref{eq:sp-order-lower-tail}.  Together,
\begin{equation}\label{eq:sp-right-laplace-extreme}
 \E e^{-p\theta s}\le C(e^{-dx}+e^{-da^\gamma}).
\end{equation}
On $\mathcal M_R^c$,
\[
 a>h,\qquad x>a^{\beta\nu_0/(\beta+\nu_0)},\qquad x/h\le a^\beta.
\]
Choosing
$\nu<\min\{\beta\nu_0/(\beta+\nu_0),\gamma\}$, polynomial absorption
in \eqref{eq:sp-insertion-moment-reduction} gives $Ce^{-ca^\nu}$.

Similarly, on $\mathcal M_L^c$,
\[
 h>a,\qquad x>h^{\alpha\nu_0/(\alpha+\nu_0)},\qquad x/h\le h^\alpha.
\]
Equations~\eqref{eq:sp-insertion-moment-reduction}
and~\eqref{eq:sp-order-lower-tail} give $Ce^{-ch^\nu}$ for
$\nu<\min\{\alpha\nu_0/(\alpha+\nu_0),1\}$.  Constants can be chosen
uniformly for $1\le p\le p_0$, proving the insertion bound.

For the restricted sums, keep these kernels fixed.  The right exceptional
support satisfies $h<a^{\beta/(\beta+\nu_0)}<a$.  Its row sum is
therefore at most $ae^{-ca^\nu}$, and its restricted column sum is at
most $\sum_{a\ge A}e^{-ca^\nu}$.  Both give the required
stretched-exponential tail after weakening the tail constant.  The left
case follows with the roles interchanged.

Finally, on the extreme sets, the displayed lower bounds for $x$ and
the inequality $xe^{-cx}\le Ce^{-cx/2}$ dominate the restricted typical
kernels by exceptional kernels with positive stretched-exponential
exponents.  The same summation argument applies.  Taking common smaller
constants and exponents in the tail bounds completes the proof.
\end{proof}

\begin{lemma}[Global domination and excursions]\label{lem:sp-domination}
For every fixed $r$ there is a deterministic nonnegative matrix $M_n(i,j)$
such that every cylinder of at most $r$ distinct marked labels and distinct
demanded ranks satisfies
\begin{equation}\label{eq:sp-global-cylinder-bound}
 \Pp(R_{n,i_e}=j_e\text{ for all }e)\le C_r\prod_eM_n(i_e,j_e).
\end{equation}
It can be chosen to have
\begin{equation}\label{eq:sp-global-row-bound}
 \sup_i\sum_jM_n(i,j)\le C_r,
\end{equation}
and, with $m=d_\sigma(j)$ at an active corner,
\begin{equation}\label{eq:sp-global-target-bound}
 \sup_iM_n(i,j)\le C_r/m
       \quad(m\le\delta n),\qquad
 \sup_iM_n(i,j)\le C_r/n
       \quad(\delta n\le j\le(1-\delta)n).
\end{equation}
Its column sums are bounded for active-corner and middle targets.  For
source and target in the same active corner it is bounded by
$C_r(H_\sigma+E_\sigma)$, apart from a fixed number of endpoint indices,
which obey the bounds in the proof below.  For some $\kappa>0$ it also has
\begin{align}
 \sup_i\sum_jM_n(i,j)
       \left(\frac{d_R(j)}{d_R(i)}\right)^\kappa&\le C_r
          &&\text{if $R$ is active},\label{eq:sp-right-weighted-row}\\
 \sup_i\sum_jM_n(i,j)
       \left(\frac{d_L(i)}{d_L(j)}\right)^\kappa&\le C_r
          &&\text{if $L$ is active}.\label{eq:sp-left-weighted-row}
\end{align}
In particular, for fixed $\ell$, the expected number of $\ell$-cycles
whose largest vertex is $\iota_\sigma(m)$ is at most $C_\ell/m$ at an
active corner and at most $C_\ell/n$ in the middle.  Among cycles rooted at
active-corner distances $A\le m\le B\le\delta n$, requiring every vertex
to have its distance in $[A,B]$ discards only $O_\ell(1)$ in expectation.
\end{lemma}

\begin{proof}
Put $p=r+1$.  Define $M_n(i,j)$ as the maximum of $\|X_{i,q}^D\|_p$
over deletion sets of size at most $p$ and valid background gaps satisfying
$|q-(j-1)|\le p+1$.  Here the gap integral defines $X_{i,q}^D$ for
every source label, whether or not it belongs to $D$.  This is a finite
maximum, and $0\le M_n(i,j)\le1$.  All estimates below are uniform in
these choices, which may vary from entry to entry.

For a cylinder with $s\le r$ marks, the ordered-simplex bound and
H\"older's inequality give
\[
 \Pp(R_{n,i_e}=j_e\text{ for all }e)
 \le \E\prod_eX_{i_e,q_e}^D
 \le\prod_e\|X_{i_e,q_e}^D\|_s
 \le\prod_eM_n(i_e,j_e).
\]
Thus it remains to establish the matrix estimates.

For targets away from a fixed number of endpoint indices, the moment
reduction in Lemma~\ref{lem:sp-insertion-moments} gives the ordinary bound
\begin{equation}\label{eq:sp-global-ordinary-kernel}
 \frac{C}{d_j}\theta_{n,i}\tau_j e^{-c\theta_{n,i}\tau_j},
 \qquad
 (d_j,\tau_j)=
 \begin{cases}
 (j,(j/n)^\alpha),&\text{active left targets},\\
 (n-j+1,(n/(n-j+1))^\beta),&\text{active right targets},\\
 (n,1),&\text{middle targets}.
 \end{cases}
\end{equation}
For sources and targets in the same active corner,
Lemma~\ref{lem:sp-insertion-moments} supplies the exceptional term
$CE_\sigma$.  Outside the right source block, the rates are bounded
below; its alternative-time argument, now at comparison depth
$n^{\beta/(\beta+1)}$, supplies an exceptional term $Ce^{-cn^\nu}$.
Outside the left source block the rates are bounded above, so
$\theta_i(j/n)^\alpha<1$ after shrinking the target block and the
moderate estimate suffices.  Middle targets have remainder $Ce^{-cn}$.
Polynomial rate prefactors are absorbed into these exponentials.

For the finitely many earliest active-left targets, the remaining-rate
floor gives $C\theta_i/n^\alpha$.  For the finitely many active-right
terminal targets, including the infinite final gap, bound the insertion
factor by survival to an earlier fixed terminal depth.  The
alternative-time estimate gives $Ce^{-ca^\nu}$ for right-source depth
$a$, and $Ce^{-cn^\nu}$ outside that block.

At an inactive left endpoint, rates are bounded above and the remaining
rate is at least $cn$; retain the bound $C\theta_i/n$.  At an inactive
right endpoint, rates are bounded below, and a linear number are also
bounded above by a fixed $M$.  At time $c_0\log(n/m)$, this latter
group has at least $c\sqrt{nm}$ expected survivors when $c_0M<1/2$.
Chernoff and the moment reduction give
\begin{equation}\label{eq:sp-inactive-terminal-kernel}
 M_n(i,j)\le
 \frac{C\theta_i}{m}(m/n)^{c\theta_i}
 +Ce^{-c\sqrt{nm}},\qquad m=n-j+1.
\end{equation}
For bounded $m$, the survival factor alone gives the same bound after
reducing $c$.  These estimates cover all gaps.

The ordinary row sums are bounded by sum--integral comparison, using
$u=\theta_i\tau_j$.  The active-corner integrals reduce to constant
multiples of $\int_0^\infty e^{-cu}\,du$; the inactive terminal integral is
\[
 \int_0^1\theta_i u^{c\theta_i-1}\,du=1/c.
\]
The middle and inactive-left row sums are bounded directly.  Exceptional
sums are bounded by Lemma~\ref{lem:sp-insertion-moments} and
stretched-exponential decay.

For active columns, the same comparison gives
\[
 \sum_a(h/a)^\alpha e^{-c(h/a)^\alpha}\le Ch,\qquad
 \sum_a(a/h)^\beta e^{-c(a/h)^\beta}\le Ch.
\]
Sources outside the corresponding corner contribute at most
$Cn(h/n)^\alpha\le Ch$ on the left and
$Cn e^{-c(n/h)^\beta}\le Ch$ on the right.  Division by the target
depth proves the column bounds.  Middle columns follow from boundedness
of $\theta e^{-c\theta}$.  The fixed endpoint columns use
$\sum_a a^{-\alpha}<\infty$ and $\sum_a e^{-ca^\nu}<\infty$.
Finally, boundedness of $ue^{-cu}$, together with the exceptional support
restrictions, gives the target maxima.  Enlarging constants handles the
finitely many small array sizes.

Choose a common $\kappa>0$ with $\kappa<\beta$ at an active right
endpoint and $\kappa<\alpha$ at an active left endpoint.  The local
ordinary row integrals remain bounded after inserting the respective
weights $(h/a)^\kappa$ and $(a/h)^\kappa$.  Local exceptional edges
have weight at most one.

For a right source of depth $a$, the total mass outside the right target
block is $O(\theta_i)=O((a/n)^\beta)$, which absorbs $(n/a)^\kappa$.
Sources outside that source block have bounded right-depth ratios, so
their ordinary row bound suffices.  For the left weight, targets outside
the left block have bounded weight.  Sources outside the left source
block have bounded rates, and their left-target contribution is bounded by
\[
 C\theta_i n^{-\alpha}
 \sum_{h\le\delta n}h^{\alpha-1}(n/h)^\kappa\le C.
\]
The fixed endpoint estimates above also satisfy these weighted bounds.
This proves both global weighted row estimates.

For cycles of length $\ell$, use the construction with a sufficient mark
budget.  At a root $v$, bound the closing edge by $C/d_\sigma(v)$,
or $C/n$ in the middle, and sum the remaining edges using the row bound.
Dropping distinctness and maximality only enlarges the sum, proving the
rooted-cycle estimates.

Finally, a largest-label right root of depth $m$ can leave $[A,B]$
only through depths greater than $B$; a left root can leave only through
depths below $A$.  Iterating the weighted row bounds along an open path,
and summing over its possible escape positions, gives respective path bounds
\[
 C_\ell(m/B)^\kappa,\qquad C_\ell(A/m)^\kappa.
\]
Multiplying by the closing-edge bound $C_\ell/m$ and summing over
$A\le m\le B$ gives $O_\ell(1)$ in either case.  The row bounds are
global, so these estimates include paths through the middle or the
opposite corner.
\end{proof}

\paragraph{Order-time concentration with a marked background.}
Let $t_h^\sigma$ denote the deterministic full-array mean quantile:
$nH_n(t_h^R)=h$ or $nG_n(t_h^L)=h$, and let
$\mu_h^\sigma=nD_n(t_h^\sigma)$.  These equations also specify the
notation when a bounded deletion is made; the deletion is treated as an
error, not as a change of quantile.
Lemma~\ref{lem:sp-quantiles} and Chernoff give, at a gap of corner
distance $h+O_r(1)$,
\begin{equation}\label{eq:sp-joint-population-concentration}
 \Pp\left(
   |T_q^D/t_h^\sigma-1|>Ch^{-1/4}
   \text{ or }|W_q^D/\mu_h^\sigma-1|>Ch^{-1/4}
        \right)
 \le Ce^{-c\sqrt h}.
\end{equation}
Here is the weight argument explicitly.  For independent Bernoulli
variables $B_i$ with means $p_i=e^{-\theta_i s}$, the inequalities
$\log(1+z)\le z$ and
$e^v-1-v\le\tfrac12v^2e^{|v|}$ give
\[
 \log\E\exp\left\{\lambda\sum_i\theta_i(B_i-p_i)\right\}
 \le\frac{\lambda^2}{2}\sum_i\theta_i^2
                         e^{-\theta_i(s-|\lambda|)}
 \le\frac{\lambda^2}{2}nD_{n,2}(s/2),
 \qquad |\lambda|\le s/2.
\]
Here $D_{n,2}(s)=n^{-1}\sum_i\theta_i^2e^{-s\theta_i}$.
At $s=t_\pm=(1\pm h^{-1/4})t_h^\sigma$, the deterministic bound is
$nD_{n,2}(s/2)\le Ch/(t_h^\sigma)^2$.
Taking $\lambda$ equal to a sufficiently small constant times
$h^{-1/4}t_h^\sigma$, with either sign, proves an exponential bound
$2e^{-c\sqrt h}$ at deviation $h^{3/4}/t_h^\sigma$.
Deleting a fixed number of clocks changes a weighted mean by at most
$C_r/t_h^\sigma$, since $\theta e^{-\theta t}\le C/t$.
The count-mean separation places $T_q^D$ between $t_-$ and $t_+$ with
error $Ce^{-c\sqrt h}$; monotonicity of the surviving weight in time and
the deterministic mean stability then prove the display.  Monotonicity of
the profile is not involved.

\begin{proposition}[Joint local insertion law]\label{prop:sp-joint-insertion}
Fix $r$.  For each of at most $r$ marked edges choose an active corner
$\sigma_e$, a source $i_e=\iota_{\sigma_e}(a_e)$ and a demanded full rank
$j_e=\iota_{\sigma_e}(h_e)$, with distinct sources and distinct ranks.
Different edges may belong to different active corners.  Suppose that
$h_0\le a_e,h_e\le\delta n$, that all edges are moderate in the sense of
Lemma~\ref{lem:sp-insertion-moments}, and that $|j_e-j_f|>2r$ for $e\ne f$.
There are $\kappa,C>0$, depending on $r$, such that
\begin{align}
 &\left|\Pp(R_{n,i_e}=j_e\text{ for all }e)
                   -\prod_eK_{\sigma_e}(a_e,h_e)\right|\notag\\
 &\qquad\le C\prod_eH_{\sigma_e}(a_e,h_e)
   \sum_e\left\{(a_e\wedge h_e)^{-\kappa}
                   +((a_e\vee h_e)/n)^\kappa\right\}.
 \label{eq:sp-joint-insertion-error}
\end{align}
For all other configurations the product domination in
\eqref{eq:sp-global-cylinder-bound} remains valid, including adjacent
demanded ranks.  In particular, if every endpoint distance belongs to
$[A,B]$, the error on the moderate separated configurations is bounded by
$C(A^{-\kappa}+(B/n)^\kappa)\prod_eH_{\sigma_e}(a_e,h_e)$.
\end{proposition}

\begin{proof}
Delete the marked labels and order the demanded ranks as in
\eqref{eq:sp-demanded-gap}.  Their separation makes the corresponding
background gaps distinct and, after increasing $h_0$, nonfinal.  Write
\[
 X_e=X_{i_e,q_e}^D,\qquad
 \xi_e=W_e(T_{q_e+1}^D-T_{q_e}^D),\qquad
 K_e=K_{\sigma_e}(a_e,h_e),\quad H_e=H_{\sigma_e}(a_e,h_e).
\]
The exact insertion identity gives
\[
 \Pp(R_{n,i_e}=j_e\text{ for all }e)=\E\prod_eX_e.
\]
Conditional on the background elimination chain, the selected normalized
gaps are independent standard exponentials.  Consequently, with $s$ the
number of marks and $Y=\prod_e\xi_e$,
\[
 \E Y=1,\qquad \E(\xi_eY)=2,\qquad \E Y^2=2^s.
\]
We use these identities without asserting independence of the insertion
factors.

The quantile and remaining-weight estimates give measurable events $G_e$
with
\[
 \Pp(G_e^c)\le Ch_e^{-1/2},
\]
on which the gap starting time and remaining rate differ relatively from
their deterministic values by $O(h_e^{-1/4})$.  Together with the sampled
endpoint expansions, these estimates imply, for some $\rho>0$,
\[
 K_e\le CH_e,\qquad X_e\le CH_e\xi_e,\qquad
 |X_e-K_e\xi_e|
 \le CH_e\left(\varepsilon_e\xi_e+\frac{\xi_e^2}{h_e}\right),
\]
where
\[
 \varepsilon_e=h_e^{-1/4}+(a_e\wedge h_e)^{-1}
 +\left(\frac{a_e\vee h_e}{n}\right)^\rho.
\]
Indeed, expand $1-e^{-\theta_e\Delta_e}$ to first order, using
$\Delta_e=\xi_e/W_e$, and compare the coefficient
$e^{-\theta_eT_{q_e}^D}\theta_e/W_e$ with $K_e$.
The exponential slack in $H_e$ absorbs the resulting powers of
$x_{\sigma_e}(a_e,h_e)$.  All estimates are uniform in the bounded
deletions and rank displacements.

Put $G=\bigcap_eG_e$.  Telescoping the product on $G$ gives
\[
 \left|\prod_eX_e-\left(\prod_eK_e\right)Y\right|
 \le C_r\left(\prod_eH_e\right)
 Y\sum_e\left(\varepsilon_e+\frac{\xi_e}{h_e}\right).
\]
Integration and the displayed gap moments bound this contribution by
$C_r\prod_eH_e\sum_e(\varepsilon_e+h_e^{-1})$.

On $G^c$, H\"older and the moderate-case estimate of
Lemma~\ref{lem:sp-insertion-moments} give
\begin{equation}\label{eq:sp-holder-weighted-bad}
 \E\left[\one_{G^c}\prod_eX_e\right]
 \le \Pp(G^c)^{1/2}\prod_e\|X_e\|_{2s}
 \le C_r\left(\prod_eH_e\right)\Pp(G^c)^{1/2}.
\end{equation}
The comparison product has the same bound by Cauchy--Schwarz, since
$\E Y^2=2^s$ and $K_e\le CH_e$.  Moreover,
\[
 \Pp(G^c)^{1/2}
 \le C_r\left(\sum_eh_e^{-1/2}\right)^{1/2}
 \le C_r\sum_eh_e^{-1/4}.
\]
Thus no independence between $G$ and the selected gaps is needed.

Choose a common $0<\kappa\le\min\{\rho,1/8\}$, and take the moderation
exponent to be $1/16$.  Every error above is bounded by a constant times
\[
 \sum_e\left\{(a_e\wedge h_e)^{-\kappa}
 +\left(\frac{a_e\vee h_e}{n}\right)^\kappa\right\}.
\]
Since $\E Y=1$, this proves \eqref{eq:sp-joint-insertion-error}.
Restricting all distances to $[A,B]$ gives the interval estimate,
absorbing the number of marks into the constant.  Global domination for
the remaining configurations follows from Lemma~\ref{lem:sp-domination}.
The empty configuration is immediate.
\end{proof}


\begin{lemma}[An integrable density envelope]\label{lem:sp-density-envelope}
Let $E_1,\ldots,E_n$ be independent continuous random variables with
densities $g_1,\ldots,g_n$, and let $R_i$ be the rank of $E_i$ in increasing
order.  Suppose $S\subseteq[n]$ and $h:\R\to[0,\infty)$ satisfy
$g_i(t)\le h(t)$ for $i\in S$ and $\int_\R h(t)\,dt<\infty$.
For every fixed $\ell\ge1$, there is a constant $K_\ell$, independent of
$n$, the clock laws, and $S$, such that
\[
 \E\#\{i\in S:i\text{ belongs to an }\ell\text{-cycle of }R\}
 \le K_\ell\int_\R h(t)\,dt.
\]
\end{lemma}

\begin{proof}
We prove the assertion with $K_\ell=3^\ell$.  For a background vector
$x=(x_1,\ldots,x_n)$, let $x^{i,t}$ be obtained by replacing $x_i$ with
$t$, and write
\[
 F_i(x)=\one\{i\text{ belongs to an }\ell\text{-cycle of the rank
 permutation of }x\}.
\]
Independence and integration in the $i$-th coordinate give
\[
 \E F_i(E)=\int_{\mathbb R}g_i(t)\,\E F_i(E^{i,t})\,dt.
\]
Thus density domination and Tonelli's theorem yield
\[
 \E\sum_{i\in S}F_i(E)
 \le\int_{\mathbb R}h(t)\,
       \E\sum_{i\in S}F_i(E^{i,t})\,dt.
\]

It remains to bound the sum inside the expectation.  Fix distinct
background values $x_1,\ldots,x_n$ and a time $t$ different from all of
them.  Let $R$ be the background rank permutation and
$N(t)=\#\{j:x_j<t\}$.  Replacing one coordinate changes the rank of
every unchanged clock by at most one; the replacement clock itself has
rank within one of $N(t)+1$.

Suppose $i$ belongs to an $\ell$-cycle of the modified rank permutation
$R^{i,t}$.  Following that cycle from the successor of $i$ back to $i$
gives a sequence
\[
 v_1=R^{i,t}(i),\quad v_2,\ldots,v_\ell=i
\]
satisfying
\[
 |v_1-(N(t)+1)|\le1,\qquad
 |v_{a+1}-R(v_a)|\le1\quad(1\le a<\ell).
\]
Indeed, exact cycle length ensures that the intermediate sources $v_a$
differ from the replaced label $i$.  There are at most three choices
for $v_1$, and at most three for each subsequent label.  Hence there
are at most $3^\ell$ such sequences.  Their final labels determine
their roots, so
\[
 \sum_{i\in S}F_i(x^{i,t})\le3^\ell.
\]

For each fixed $t$, independence and the existence of densities ensure
that the background values are almost surely distinct and avoid $t$.
Applying this deterministic bound in the preceding integral proves
\[
 \E\#\{i\in S:i\text{ belongs to an }\ell\text{-cycle of }R\}
 \le3^\ell\int_{\mathbb R}h(t)\,dt.
\]
\end{proof}


\subsection{Cycle bounds and logarithmic traces}
\label{subsec:sp-cycle-traces}

Throughout this subsection and the next two we assume
Assumption~\ref{ass:simple-power-profile}.  Constants may depend on the
profile and on a fixed upper bound for the number of marked clocks.  Write
\[
 d_R(i)=n+1-i,\qquad d_L(i)=i.
\]
A cycle is always rooted at its largest label, as in the introduction.
Its root therefore has the smallest right distance and the largest left
distance among its vertices.

We first record consequences of the insertion estimates that hold before
any asymptotic approximation is made.  The domination matrix in
Lemma~\ref{lem:sp-domination} can be chosen for any fixed number of marked
clocks.  We denote it by $F$; thus the probability of a prescribed finite
collection of assignments is bounded by the product of its $F$ entries.
Its row sums are bounded.  For an active target corner $\sigma$, its
entries satisfy
\begin{equation}\label{eq:sp-F-target}
 F(i,j)\le \frac{C}{d_\sigma(j)}
 \quad\text{when }j\text{ belongs to a sufficiently small fixed
 neighborhood of that corner}.
\end{equation}
The same lemma supplies the weighted row bounds, with some $\kappa>0$,
\begin{align}
 \sum_jF(i,j)\left(\frac{d_R(j)}{d_R(i)}\right)^\kappa&\le C,
 &&\text{at an active right corner},\label{eq:sp-row-right}\\
 \sum_jF(i,j)\left(\frac{d_L(i)}{d_L(j)}\right)^\kappa&\le C,
 &&\text{at an active left corner}.\label{eq:sp-row-left}
\end{align}
Here and below the global versions of these bounds include paths that
leave the endpoint neighborhood.  Products along any fixed number of
steps have the same bounds with a different constant, because the weight
in each display telescopes along the path.

\begin{lemma}[Root and excursion bounds]\label{lem:sp-root-excursion}
Fix $L$.  At an active corner $\sigma$, the expected number of
$\ell$-cycles with root distance $m$ is at most $C_L/m$, uniformly for
$\ell\le L$ and $m$ in a sufficiently small fixed endpoint neighborhood.
The contribution of a fixed middle interval, and that of every inactive
endpoint, is bounded uniformly in $n$.

For $1\le A\le B$ within an active endpoint neighborhood, the expected
number of cycles whose root distance lies in $[A,B]$ but which have a
vertex outside that distance interval is $O_L(1)$.  More generally, for
$R\ge1$, the expected number with root distance in $[A,B]$ and with some
vertex at logarithmic distance at least $\log R$ from the root is at most
\begin{equation}\label{eq:sp-log-excursion}
 C_L(1+\log(B/A))R^{-\kappa_L}.
\end{equation}
\end{lemma}

\begin{proof}
The active-corner root bound and the $O_\ell(1)$ interval-exit bound follow
from Lemma~\ref{lem:sp-domination}.  Taking maxima of the constants and a
common smaller endpoint neighborhood makes both estimates uniform for
$1\le\ell\le L$.

In a fixed middle interval or a sufficiently small neighborhood of an
inactive endpoint, the sampled rates lie in some fixed interval
$[d,M]\subset(0,\infty)$.  Their densities are dominated by
\[
 h(t)=Me^{-dt}\one_{\{t\ge0\}},\qquad
 \int_{\mathbb R}h(t)\,dt=M/d.
\]
Lemma~\ref{lem:sp-density-envelope} therefore bounds the expected number
of vertices in these regions belonging to an $\ell$-cycle by
$3^\ell M/d$.  This also bounds the number of cycles rooted there,
uniformly for $\ell\le L$.

It remains to prove the logarithmic-excursion estimate.  Fix a root of
endpoint distance $m$.  For $1\le R\le2$, its contribution is at most
$C_L/m$, which is bounded by $C'_Lm^{-1}R^{-\kappa_L}$ after increasing
the constant.

Suppose $R>2$.  At the right corner, the largest-label root has minimum
endpoint distance.  Thus logarithmic distance at least $\log R$ requires
a vertex of distance at least $Rm$, and hence greater than $Rm/2$.
Apply the global weighted row bound along the open path to that vertex,
and the ordinary row bound thereafter.  Bounding the closing edge by
$C_L/m$ gives
\[
 \frac{C_L}{m}\left(\frac{m}{Rm/2}\right)^{\kappa_L}
 \le \frac{C'_L}{m}R^{-\kappa_L}.
\]
At the left corner, the root has maximum endpoint distance.  The
corresponding event requires a vertex of distance at most $m/R$, hence
below $2m/R$.  The left weighted row bound gives the same estimate.
Summing over the possible vertex positions changes only the constant;
the global bounds include paths through the middle or the opposite corner.

Finally,
\[
 \sum_{\substack{m\in\mathbb N\\A\le m\le B}}\frac1m
 \le1+\log(B/A).
\]
Summing the per-root estimate proves \eqref{eq:sp-log-excursion}.
A common positive exponent and common constants may be chosen over the
finitely many cycle lengths and active corners.
\end{proof}

For an active corner put
\[
 K_R(a,b)=\frac1b\,\beta q_R(a/b)^\beta
              e^{-q_R(a/b)^\beta},\qquad
 K_L(a,b)=\frac1b\,\alpha q_L(a/b)^{-\alpha}
              e^{-q_L(a/b)^{-\alpha}}.
\]
These are deterministic ideal kernels on endpoint distances.  They are
not asserted to be globally uniform pointwise approximations to an
averaged insertion probability.  They are used below only through
Proposition~\ref{prop:sp-joint-insertion} and its domination estimates.
In logarithmic coordinates their increment densities are, respectively,
$p_{\beta,q_R}$ and the reflection of $p_{\alpha,q_L}$.  Denote the
appropriate density by $p_\sigma$.  Reflection does not change its return
density at zero, and
\[
 b_{\sigma,\ell}=\frac1\ell p_\sigma^{*\ell}(0).
\]

For integer cutoffs $A\le B$, let $\lambda_{n,\sigma,\ell}(A,B)$ be the
sum of $\prod_{j=1}^\ell K_\sigma(a_j,a_{j+1})$, $a_{\ell+1}=a_1$,
over distinct vertices in $[A,B]$, divided by $\ell$.  If a root interval
is specified, retain only cycles whose largest-label root has its
distance in that interval.  These exact ideal sums, including their
boundary corrections, will be the means used in the factorial-moment
comparison.

\begin{lemma}[Ideal traces]\label{lem:sp-ideal-traces}
For fixed $\ell$, as $A\to\infty$, uniformly for $B\ge A$,
\begin{equation}\label{eq:sp-ideal-total}
 \lambda_{n,\sigma,\ell}(A,B)
 =b_{\sigma,\ell}\log(B/A)+O_\ell(1)
   +O_\ell\bigl((1+\log(B/A))/A\bigr).
\end{equation}
Let $A_n=\lceil\exp((\log n)^{1/4})\rceil$ and
$B_n=\lfloor n/A_n\rfloor$.  For a fixed interval
$I=(a,b]$ with $0<a<b<1$, the ideal sum restricted to vertices in
$[A_n,B_n]$ and root distance in $(n^a,n^b]$ satisfies
\begin{equation}\label{eq:sp-ideal-spatial}
 \lambda_{n,\sigma,\ell,I}
 =b_{\sigma,\ell}|I|\log n
   +O_\ell\bigl((1+\log n)^K A_n^{-\kappa}\bigr)
\end{equation}
for some $\kappa>0$ and finite $K$.
\end{lemma}

\begin{proof}
Write $\lambda_\sigma(A,B)=\lambda_{n,\sigma,\ell}(A,B)$ and
$p=p_\sigma$.  For $y_0=y_\ell=0$, put
\[
 W(\boldsymbol y)=\prod_{j=0}^{\ell-1}p(y_j-y_{j+1}),
 \qquad
 R(\boldsymbol y)=\max_{0\le j<\ell}y_j-\min_{0\le j<\ell}y_j.
\]
The explicit density satisfies $p(x)+|p'(x)|\le Ce^{-g|x|}$.
Consequently $\int e^{cR}W<\infty$ for some $c>0$, by bounding one
factor and integrating the remaining increments.  In logarithmic
coordinates, the continuous trace over an interval of width $T$ is
\[
 \frac1\ell\int_{\mathbb R^{\ell-1}}(T-R)_+W
 =b_{\sigma,\ell}T-d_{\sigma,\ell}+O_\ell(e^{-cT}),
 \qquad
 d_{\sigma,\ell}=\frac1\ell\int_{\mathbb R^{\ell-1}}RW.
\]
Indeed, $\int W=p^{*\ell}(0)$, and
$(T-R)_+=T-R+(R-T)_+$; exponential integrability bounds the last
term.  For $\ell=1$, the range and boundary constant are zero.

The identity
\[
 K_\sigma(u,v)=\frac{p(\log u-\log v)}{v}
\]
allows comparison with logarithmic cells
$[\log m,\log(m+1))$.  Their widths are
$m^{-1}(1+O(A^{-1}))$.  Telescoping the product and using the derivative
envelope gives total error $O_\ell((1+T)/A)$, including the upper
boundary adjustment, where $T=\log(B/A)$.  Removing repeated vertices
costs $O_\ell(A^{-1})$: bounded kernel rows and the entry bound $C/v$
bound each prescribed coincidence by
$C_\ell\sum_{m=A}^B m^{-2}$.
Thus, uniformly for $B\ge A$,
\[
 \lambda_\sigma(A,B)
 =b_{\sigma,\ell}\log(B/A)-d_{\sigma,\ell}
 +O_\ell\!\left(\frac{1+\log(B/A)}{A}
                  +e^{-c\log(B/A)}\right).
\]
This implies~\eqref{eq:sp-ideal-total}.

For the spatial estimate, set $P=\lfloor n^a\rfloor$ and
$Q=\lfloor n^b\rfloor$.  For large $n$, $A_n\le P<Q<B_n$.
Since the canonical root has maximum left depth and minimum right
depth, respectively, the exact identities are
\[
 \lambda_{n,\sigma,\ell,I}
 =
 \begin{cases}
  \lambda_L(A_n,Q)-\lambda_L(A_n,P),&\sigma=L,\\
  \lambda_R(P+1,B_n)-\lambda_R(Q+1,B_n),&\sigma=R.
 \end{cases}
\]
Apply the refined expansion to both terms.  The constants
$d_{\sigma,\ell}$ cancel.  Every interval appearing here has
logarithmic width at least $c_I\log n$, so its exponential remainder
is $O(A_n^{-1})$, since $\log A_n=o(\log n)$.  The lattice errors are
$O_\ell((1+\log n)/A_n)$, and rounding changes the difference of
logarithmic widths from $(b-a)\log n$ by $O(n^{-a})$.  Hence
\[
 \lambda_{n,\sigma,\ell,I}
 =b_{\sigma,\ell}(b-a)\log n
 +O_\ell\!\left(\frac{1+\log n}{A_n}\right),
\]
which proves~\eqref{eq:sp-ideal-spatial} with $K=\kappa=1$.
\end{proof}

\subsection{Mixed factorial moments}
\label{subsec:sp-factorial-moments}

We now compare actual core cycles with the ideal sums.  All constants in
this comparison may depend on the fixed factorial-moment order.  This
dependence is harmless; no estimate uniform in that order is needed.

\begin{lemma}[Overlaps and nearby ranks]\label{lem:sp-collisions}
Fix a finite number $s$ of proposed cycles of lengths at most $L$, all
with vertices in endpoint cores of distances $[A,B]$.  Suppose their
joint weight is bounded by a product of a matrix having bounded row sums
and entries at most $C/b$ for a core target of distance $b$.  The sum of
such weights over assignments in which the actual labels in two specified
vertex slots differ by at most a fixed $D$ is at most
\begin{equation}\label{eq:sp-collision-sum}
 C_{s,L,D}A^{-1}(1+\log(B/A))^{s-1}.
\end{equation}
For slots in the same corner this is equivalent to their endpoint
distances differing by at most $D$.  This includes equal vertices.
The same bound, with a changed constant,
holds for the union over every pair of slots.
\end{lemma}

\begin{proof}
Put $H=1+\log(B/A)$, and write $d_\sigma(i)$ for the endpoint
distance of label $i$.  Restrict $F$ to each corner core, obtaining
$F_\sigma$.  Bounded rows and the target bound imply, for every
fixed $j\ge1$,
\[
 (F_\sigma^j)(u,v)\le \frac{C_j}{d_\sigma(v)},
 \qquad
 \operatorname{Tr}(F_\sigma^j)\le C_jH.
\]
For two distinct slots in one cycle, split the cycle into the two
positive-length paths between them.  For slots in different cycles,
root each cycle at its selected slot.  In either case, summing the
other vertices of the affected cycles gives at most
\[
 \frac{C_{s,L}}{d_\sigma(u)d_\tau(v)}.
\]

Each label has at most $2D+1$ labels within distance $D$.  Since
endpoint distance is injective on each core and
$2/(xy)\le x^{-2}+y^{-2}$,
\[
 \sum_{\substack{u\text{ in the }\sigma\text{ core},\
                  v\text{ in the }\tau\text{ core}\\ |u-v|\le D}}
 \frac1{d_\sigma(u)d_\tau(v)}
 \le (2D+1)\sum_{m=A}^B m^{-2}
 \le \frac{2(2D+1)}A.
\]
This includes mixed corners without requiring their cores to be
disjoint.  The unaffected cycles contribute at most
$C_{s,L}H^{s-1}$, since $H\ge1$, proving the specified-pair bound.
Summing over at most $(sL)^2$ slot pairs proves the union bound.
For slots in the same corner, label differences equal
endpoint-distance differences.  The argument also permits different
matrices on different cycles when their bounds are uniform.
\end{proof}

Consider any fixed finite list of categories $c$, each specifying a
corner, a length, and either its whole core or a prescribed root interval.
For categories of the same corner and length require the root intervals
to be disjoint.  Let $Y_{n,c}$ count the actual cycles in that category
whose every vertex lies in its core, and let $\lambda_{n,c}$ be the exact
ideal sum with precisely these restrictions.

\begin{proposition}[Core factorial moments]\label{prop:sp-factorial-moments}
With $A_n,B_n$ as in Lemma~\ref{lem:sp-ideal-traces}, for every fixed
multi-index $\boldsymbol r=(r_c)$ there are
$\kappa_{\boldsymbol r}>0$ and $K_{\boldsymbol r}<\infty$ such that
\begin{equation}\label{eq:sp-factorial-comparison}
 \left|\E\prod_c(Y_{n,c})_{r_c}
             -\prod_c\lambda_{n,c}^{r_c}\right|
 \le C_{\boldsymbol r}(1+\log n)^{K_{\boldsymbol r}}
                    A_n^{-\kappa_{\boldsymbol r}}.
\end{equation}
Here $(x)_r=x(x-1)\cdots(x-r+1)$.
\end{proposition}

\begin{proof}
Let $s=\sum_c r_c$, $m=\sum_c r_c\ell_c$, and
$\mathcal D=\prod_c\ell_c^{r_c}$.  The case $s=0$ is immediate.
Write $A=A_n$, $B=B_n$, and $H=1+\log(B/A)$.

Expand the mixed falling factorial into ordered lists of cycles and
represent each cycle by an ordered vertex tuple.  For large $n$,
the endpoint cores are disjoint; together with category disjointness,
this ensures that all selected permutation cycles are distinct and
hence vertex-disjoint.  Thus the actual moment is the sum of joint
rank probabilities over globally distinct assignments, divided by
$\mathcal D$.  The product of exact ideal means has the same
representation with ideal edge weights and distinctness required
only within each cycle.  Both sums retain the prescribed root
restrictions.

First discard assignments with two labels within distance $2m$.
Lemma~\ref{lem:sp-collisions}, applied to the actual domination matrix
and to the ideal kernels, bounds their contribution by
$CA^{-1}H^{s-1}$.

Next discard nonmoderate edges.  In endpoint coordinates, either
matrix $F$ has bounded rows, entries at most $C/b$, and, for a
suitable $t>0$,
\[
 \sum_b F(a,b)V_\sigma(b)\le CV_\sigma(a),
 \qquad
 V_L(a)=a^{-t},\quad V_R(a)=a^t.
\]
For any specified edge $a_i\to a_{i+1}$, cut that edge, bound it by
$C/a_{i+1}$, and propagate the weighted-row estimate along the
remaining path.  This gives
\[
 \sum_{a_1,\ldots,a_\ell\in[A,B]}
 \left(\prod_{j=1}^{\ell}F(a_j,a_{j+1})\right)
 \frac{V_\sigma(a_i)}{V_\sigma(a_{i+1})}
 \le C_\ell H,
 \qquad a_{\ell+1}=a_1.
\]
Writing $\gamma_L=\alpha$ and $\gamma_R=\beta$, the displayed ratio
equals $x_\sigma(a_i,a_{i+1})^{t/\gamma_\sigma}$.
Failure of moderation therefore makes it at least $A^\eta$ for some
$\eta>0$.  Summing over edges and the remaining cycles bounds all
nonmoderate assignments by $CA^{-\eta}H^s$.  This applies to the full
domination matrix, including its exceptional terms, and to the ideal
kernels.

On the remaining assignments, apply
Proposition~\ref{prop:sp-joint-insertion} jointly to all $m$ marked
vertices, including both corners in the same race.  Its error is at
most
\[
 C\bigl(A^{-\kappa}+(B/n)^\kappa\bigr)
 \prod_e H_{\sigma_e}(a_e,h_e).
\]
The envelope kernels have bounded rows and target bounds, so their
product sums to at most $CH^s$.  Since $B/n\le A^{-1}$, this
contribution is $CA^{-\kappa}H^s$.  Category restrictions remain in
the comparison and are dropped only when bounding nonnegative
error sums.

Combining the three contributions, reducing the positive exponent
if necessary, and using $H\le1+\log n$ gives
\[
 C_{\boldsymbol r}(1+\log n)^{s+1}A_n^{-\kappa_{\boldsymbol r}}.
\]
Division by $\mathcal D\ge1$ proves~\eqref{eq:sp-factorial-comparison},
with $K_{\boldsymbol r}=s+1$.
\end{proof}

\subsection{Power-law CLTs and localization}
\label{subsec:sp-power-conclusions}

\begin{proof}[Proof of Theorem~\ref{thm:power-clt}]
Let $Y_{n,\sigma,\ell}$ denote the core counts and
$\lambda_{n,\sigma,\ell}$ their exact ideal means.  By
Lemma~\ref{lem:sp-ideal-traces},
\[
 \lambda_{n,\sigma,\ell}
 =b_{\sigma,\ell}\log n+O(1+\log A_n).
\]
Proposition~\ref{prop:sp-factorial-moments} determines their limiting
centered mixed moments: conversion from falling factorials to ordinary
powers introduces only fixed powers of $\log n$, while
\[
 (1+\log n)^K A_n^{-\kappa}\longrightarrow0
 \qquad(K<\infty,\ \kappa>0).
\]
Thus the mixed moments of the counts centered at their exact ideal
means and divided by their square roots agree, up to $o(1)$, with the
corresponding products of normalized Poisson moments.

For completeness, this comparison can be made algebraically.  Define
the linear functional $\mathcal L_\lambda$ on polynomials by
$\mathcal L_\lambda((x)_r)=\lambda^r$.  Its centered moments
$M_r(\lambda)=\mathcal L_\lambda((x-\lambda)^r)$ satisfy
\[
 M_0=1,\qquad M_1=0,\qquad
 M_{r+1}(\lambda)
 =\lambda\sum_{j=0}^{r-1}\binom{r}{j}M_j(\lambda).
\]
After normalization by $\lambda^{(r+1)/2}$, induction gives the standard
Gaussian moments as $\lambda\to\infty$.  Hence the normalized core
counts have limiting mixed moments equal to those of independent
standard normals.  Taylor expansion of the characteristic function of
each linear combination, with the remainder controlled by even moments,
gives convergence of characteristic functions by taking $n\to\infty$
first and then increasing the Taylor order.  L\'evy's theorem yields
joint Gaussian convergence.

Since $\log A_n=O((\log n)^{1/4})$, replacing the ideal means by
$b_{\sigma,\ell}\log n$ preserves this limit.
Lemma~\ref{lem:sp-root-excursion} also gives
\[
 \E\left|C_{n,\ell}-\sum_{\sigma\ {\rm active}}
 Y_{n,\sigma,\ell}\right|
 =O(1+\log A_n)=o(\sqrt{\log n}),
\]
by summing the bounds for roots outside the core, cycles leaving the
core, and the middle and inactive endpoint regions.  The normalized
error therefore vanishes in $L^1$.

For each length, the corner limits combine with weights
$\sqrt{b_{\sigma,\ell}/B_\ell}$, whose squares sum to one.  This proves
the asserted independent standard normal limits.  The first
factorial-moment comparison and the same omission bound give
$\E C_{n,\ell}\sim B_\ell\log n$.
\end{proof}

\begin{proof}[Proof of Theorem~\ref{thm:power-spatial}]
Apply Proposition~\ref{prop:sp-factorial-moments} to the categories
$(\sigma,\ell,I_j)$, retaining the largest-label root.
Lemma~\ref{lem:sp-ideal-traces} gives their exact ideal means as
\[
 \lambda_{n,\sigma,\ell,I_j}
 =b_{\sigma,\ell}|I_j|\log n+o(\sqrt{\log n}).
\]
The mixed-moment and characteristic-function argument in the preceding
proof therefore gives independent Gaussian limits for all normalized
core counts.

Since $\overline{I_j}\subset(0,1)$, every root in $n^{I_j}$ lies in
$[A_n,B_n]$ for sufficiently large $n$.
Lemma~\ref{lem:sp-root-excursion} bounds the expected number of such
cycles having a vertex outside the core by $O(1)$.  Consequently the
normalized difference between the full and core counts tends to zero
in $L^1$, proving the joint CLT.  The first factorial-moment comparison
and this same bound prove the stated mean asymptotics.

Finally, the largest-label root has minimum right distance and maximum
left distance.  If a vertex's logarithmic distance differs from the
root's by more than $\delta\log n$, the corresponding distance ratio
exceeds $n^\delta$.  The excursion bound in
Lemma~\ref{lem:sp-root-excursion} therefore bounds the expected number
of such cycles by
\[
 C_L(1+\log n)n^{-\kappa_L\delta}\longrightarrow0,
\]
which proves localization for every vertex.
\end{proof}


\subsection{The critical pole}

We first record the martingale criterion used for fixed points.
\begin{proposition}\label{prop:sp-fixed-martingale}
Let $I_{n,k}\in\{0,1\}$ be adapted to a filtration, and put
$p_{n,k}=\E(I_{n,k}\mid\cF_{n,k-1})$.  If $v_n\to\infty$ and
\[
 \sum_kp_{n,k}=v_n+o_{\Pp}(\sqrt{v_n}),
 \qquad \sum_kp_{n,k}^2=o_{\Pp}(v_n),
\]
then $(\sum_kI_{n,k}-v_n)/\sqrt{v_n}\Longrightarrow\mathcal N(0,1)$.
If additionally $\E\sum_kp_{n,k}\sim v_n$, then
$\E\sum_k I_{n,k}\sim v_n$.
\end{proposition}
\begin{proof}
Write $A_n=\sum_kp_{n,k}$.  Stop each row before its compensator would
exceed $2v_n$.  The stopped and original rows agree on
$\{A_n\le2v_n\}$, whose probability tends to one.  Thus the stopped
row inherits both convergence assumptions, and its normalized square
sum is bounded by $2$, so
\[
 \E\!\left[\frac1{v_n}\sum_kp_{n,k}^2\right]\longrightarrow0.
\]
Hereafter work with the stopped row and suppress $n$.

Fix $t\in\mathbb R$, put $u=t/\sqrt v$, $a=e^{iu}-1$, and define
\[
 S_k=\sum_{j\le k}I_j,\qquad A_k=\sum_{j\le k}p_j,
 \qquad Z_k=e^{iuS_k-aA_k}.
\]
The conditional Bernoulli identity gives
\[
 \E(Z_k\mid\cF_{k-1})
 =Z_{k-1}e^{-ap_k}(1+ap_k).
\]
Since $A_k\le2v$, we have $|Z_k|\le C_t$.  Also
$\lvert e^{-ap}(1+ap)-1\rvert\le C|a|^2p^2$ for large $n$.
Telescoping expectations therefore yields
\[
 |\E Z_N-1|
 \le C_t\,\E\!\left[\frac1v\sum_kp_k^2\right]
 \longrightarrow0.
\]

The first hypothesis implies $a(A_N-v)\to0$ in probability.  Both
$e^{aA_N-iuv}$ and $e^{av-iuv}$ have modulus at most one, so their
difference tends to zero in $L^1$.  Consequently
\[
 \E e^{iu(S_N-v)}
 =e^{av-iuv}\E Z_N+o(1)
 \longrightarrow e^{-t^2/2},
\]
because $v(e^{iu}-1-iu)\to-t^2/2$.  L\'evy's continuity theorem,
followed by removal of stopping, proves the CLT.  Finally,
$\E\sum_kI_{n,k}=\E\sum_kp_{n,k}$, proving the expectation assertion.
\end{proof}

\begin{proof}[Proof of Theorem~\ref{thm:critical-pole}]
Write $c=c_L$.  The profile assumptions give
\begin{equation}\label{eq:sp-critical-global-comparison}
 \theta_{n,i}\asymp \frac ni,
 \qquad
 \sum_{i=1}^n\left|\theta_{n,i}-\frac{cn}{i-1/2}\right|\le Cn.
\end{equation}
Indeed, the error from the pole is bounded by an integrable power
near zero and is bounded on intervals separated from zero.
Comparing with the pure-pole integrals and their Riemann sums
therefore gives
\begin{equation}\label{eq:sp-critical-transforms}
 G_n(t)=ct\log(1/t)+O(t+n^{-1}),\qquad
 D_n(t)=c\log(1/t)+O(1+(nt)^{-1}),
\end{equation}
where $G_n(t)=n^{-1}\sum_i(1-e^{-t\theta_{n,i}})$ and
$D_n(t)=n^{-1}\sum_i\theta_{n,i}e^{-t\theta_{n,i}}$.
The same rate comparison shows that deleting at most $m$ labels
leaves total rate at least
\[
 n\{c\log(n/m)-C\}.
\]
To see this, restrict the harmonic comparison to $i\ge m$: the
deleted terms have total at most $m/m=1$.

Choose a sufficiently small fixed $\eps>0$, and put
\[
 A=\lceil\sqrt n\rceil,\qquad B=\lfloor\eps n\rfloor,
 \qquad z_m=\log(n/m).
\]
Let $W_m$ be the rate remaining before draw $m$, and $J_m$ indicate
that label $m$ remains.  Then $p_{n,m}=\theta_{n,m}J_m/W_m$.

For $A\le m\le B$, test the arrival count at
\[
 t_-=\frac{m}{16cnz_m},\qquad t_+=\frac{8m}{cnz_m}.
\]
The population estimates place its means below $m/2$ and above
$2m$, respectively.  Chernoff bounds, the remaining-rate lower
bound, and the estimate for $D_n(t_-)$ give
\[
 \begin{aligned}
 W_m&\ge n(cz_m-C),\\
 \E W_m&\le cnz_m+Cn(1+\log z_m+z_m/m),\\
 \Pp(J_m=0)&\le C(z_m^{-1}+m^{-1}).
 \end{aligned}
\]
For the second estimate, bound $W_m$ by the surviving rate at $t_-$
when fewer than $m$ arrivals have occurred.  The complementary
contribution is at most $Cn^2e^{-c'm}\le Cn/m$, since $n\le m^2$.
For the third, use $\Pp(E_{n,m}\le t_+)\le\theta_{n,m}t_+$ and the
lower arrival-tail bound.

For small enough $\eps$, the first estimate gives
$W_m\ge cnz_m/2$.  Combining the three estimates with the sampled
endpoint expansion yields
\[
 \E\left|p_{n,m}-\frac1{mz_m}\right|
 \le \frac C{mz_m}
 \left\{\frac{1+\log z_m}{z_m}+\frac1m+(m/n)^\eta\right\}.
\]
Here the reciprocal error follows from
\[
 \E|W_m-cnz_m|
 \le \E W_m-cnz_m+2Cn.
\]
Integral comparison shows that the displayed probability errors
have bounded sum.

Define $q_{n,m}=1/(mz_m)$ on $[A,B]$ and zero elsewhere.  Before $A$,
the deterministic bound $p_{n,m}\le C/(mz_m)$ has bounded sum.
Beyond $B$, the rates have uniform positive lower and upper bounds,
so Lemma~\ref{lem:sp-density-envelope} gives
$\sum_{m>B}\E p_{n,m}=O(1)$.  Thus
\[
 \sum_m\E|p_{n,m}-q_{n,m}|=O(1),\qquad
 \sum_mq_{n,m}=\log\log n+O(1),\qquad
 \sum_mq_{n,m}^2=O(1).
\]
Since $0\le p_{n,m},q_{n,m}\le1$, these imply
\[
 \sum_mp_{n,m}=\log\log n+O_{\Pp}(1),\qquad
 \E\sum_mp_{n,m}=\log\log n+O(1),\qquad
 \E\sum_mp_{n,m}^2=O(1).
\]
Proposition~\ref{prop:sp-fixed-martingale} proves the fixed-point
CLT and mean asymptotic.

Now fix $\ell\ge2$.  Choose a sufficiently large fixed $M=M_\ell$,
and decrease $\eps$ if necessary.  Cycles meeting labels below $M$
number at most $M$; cycles meeting labels above $B$ have bounded
expected number by Lemma~\ref{lem:sp-density-envelope}.  It remains
to consider cycles contained in $[M,B]$.

Put $L_j=\log(n/j)$.  The deleted-rate lower bound and the arrival
estimate at a small multiple of $j/(nL_j)$, combined with the
exponential spacing moments, give an insertion domination kernel
\[
 F_n(i,j)=\frac{C_\ell}{iL_j}
 \exp\!\left(-\gamma_\ell\frac{j}{iL_j}\right),
 \qquad M\le i,j\le B.
\]
This estimate is uniform under the finitely many deletions and rank
shifts needed for an $\ell$-cycle.  The arrival-tail error is absorbed
into the displayed exponential because $j/(iL_j)\le j$.
H\"older's inequality therefore bounds each cycle probability by
the product of its $F_n$-weights.

Root the cycle at its largest label $m=i_1$, and set
\[
 z=L_m,\qquad v_j=\log(m/i_j)\ge0,\qquad
 S=\sum_{j=1}^{\ell}v_j,\qquad v_1=v_{\ell+1}=0.
\]
A closed path has some drop $v_j-v_{j+1}\ge S/(2\ell^2)$.  Hence
\[
 \sum_{j=1}^{\ell}\frac{i_{j+1}}{i_jL_{i_{j+1}}}
 \ge \frac{e^{S/(2\ell^2)}}{z+S}.
\]
The elementary bound $e^{-\gamma x}\le Cx^{-1/2}$ then gives
\[
 \exp\!\left(-\gamma_\ell
 \frac{e^{S/(2\ell^2)}}{z+S}\right)
 \le C_\ell z^{1/2}e^{-S/(8\ell^2)}.
\]
Since every $L_{i_j}\ge z$ and $\ell\ge2$, writing
$a=1/(8\ell^2)$ yields
\[
 \prod_{j=1}^{\ell}F_n(i_j,i_{j+1})
 \le \frac{C_\ell}{mz^{3/2}}
 \prod_{j=2}^{\ell}\frac1{i_j}\left(\frac{i_j}{m}\right)^a.
\]
Each remaining vertex sums to at most $1+1/a$, even after dropping
distinctness.  The contribution rooted at $m$ is therefore at most
$C_\ell/(mL_m^{3/2})$.  Finally,
\[
 \sum_{m=M}^{B}\frac1{mL_m^{3/2}}
 \le C\left(1+\int_{\log(1/\eps)}^\infty z^{-3/2}\,dz\right)<\infty.
\]
Adding the omitted cycles proves the uniform expectation bound;
finitely many smaller $n$ are absorbed into the constant.
\end{proof}


\subsection{The sampling grid and Sukhatme weights}

The midpoint convention keeps the statement of the profile assumption
simple.  The following observation connects it to the exact weights in
Corollary~\ref{cor:sukhatme}.

\begin{lemma}[Interior-grid sampling]\label{lem:sp-interior-grid}
The conclusions of Theorems~\ref{thm:power-clt}
and~\ref{thm:power-spatial} also hold if the sampling rule in
Assumption~\ref{ass:simple-power-profile} is replaced by
\[
 \theta_{n,i}=f(i/(n+1)),\qquad 1\le i\le n.
\]
The same replacement is valid in Theorem~\ref{thm:critical-pole}.
\end{lemma}

\begin{proof}
All sample points remain in $(0,1)$.  The endpoint distances are now
$m/(n+1)$ at both ends.  On an active endpoint block the deterministic
expansions used above become
\[
 \theta_{n,n-m+1}
 =c_R(m/(n+1))^\beta(1+O((m/n)^\eta)),\qquad
 \theta_{n,m}
 =c_L(m/(n+1))^{-\alpha}(1+O((m/n)^\eta)).
\]
Compared with their leading expressions on the midpoint grid, these
add relative errors $O(m^{-1}+n^{-1})$.  In the proof of
Lemma~\ref{lem:sp-populations}, the mesh is $1/(n+1)$ rather than $1/n$;
division of the sums by $n$ adds a factor $1+1/n$.  The error in an
unweighted sum remains $O(n^{-1})$, and that in a weighted sum remains
$O((nt)^{-1})$.  Thus the population and quantile estimates, including
their leading constants, are unchanged.  Endpoint comparability and
the positive upper and lower bounds off the active neighborhoods are
also unchanged.

The insertion and cycle arguments use precisely these deterministic
estimates and bounds.  The additional error on their core is
$O(A_n^{-1}+n^{-1})$, which is already allowed by the mixed factorial
moment estimates.  The bounds outside the core remain uniform.
This proves the two power-law assertions.  For the critical pole the
same sum--integral estimates give $G_n(t)=c_Lt\log(1/t)+O(t+n^{-1})$
and $D_n(t)=c_L\log(1/t)+O(1+(nt)^{-1})$ in the range used above;
global comparability with $n/i$ persists.  Hence the critical
fixed-point and longer-cycle arguments apply as well.
This is a verification of the proof inputs for the second grid; it
does not assert a coupling between the cycle counts on the two grids.
\end{proof}


\section{Independent-clock extensions}\label{sec:clock-extensions}

The bulk density~\eqref{eq:rho} uses the exponential form of the Luce race,
but the last-position estimate has a distribution-free core.  This identifies
a broader class of ranking models for which the same fixed-point strategy is
available.

Let $T_1,\ldots,T_n$ be independent continuous random variables with
densities $g_i$ and survival functions $\overline G_i$.  Order them
increasingly and write
\[
 R_i=1+\sum_{j\ne i}\one\{T_j<T_i\},\qquad
 S(t):=\sum_{i=1}^n\overline G_i(t).
\]
As before, put $k_m=n-m+1$.

\begin{proposition}
\label{prop:general-clock-tail}
Let $1\le M\le n$, choose $B$ with $B-1\ge2M$, and let $s$ satisfy
$S(s)=B$.  Suppose that an integrable function $h$ satisfies
\begin{equation}\label{eq:clock-density-envelope}
 g_{k_m}(t)\le h(t),\qquad 1\le m\le M,\quad t\ge s.
\end{equation}
Then, for a universal $c>0$,
\begin{equation}\label{eq:general-clock-tail}
 \sum_{m=1}^M\Pp(R_{k_m}=k_m)
 \le M e^{-cB}+2\int_s^\infty h(t)\,dt.
\end{equation}
\end{proposition}

\begin{proof}
Conditioning on $T_{k_m}=t$ gives the exact identity
\[
 \Pp(R_{k_m}=k_m)=\int_{-\infty}^{\infty}g_{k_m}(t)
 \Pp\left(\sum_{j\ne k_m}\one\{T_j>t\}=m-1\right)dt.
\]
For $t\le s$, the survivor count in parentheses has mean at least
$S(t)-1\ge B-1\ge2M$, so its lower tail is at most $e^{-cB}$ by a Chernoff
bound.  The integral over this range, summed over $m$, is at most
$Me^{-cB}$.

For fixed $t$, couple all the inner probabilities using one independent
sample of the clocks.  If $N(t)=\sum_j\one\{T_j>t\}$, the event in the last
display can hold only for $m=N(t)$ or $m=N(t)+1$.  Hence the sum of its
indicators over $m\le M$ is at most two.  On $t\ge s$, use
\eqref{eq:clock-density-envelope}, take expectations, and integrate to obtain
the second term in~\eqref{eq:general-clock-tail}.
\end{proof}

Suppose more generally that the clock attached to label $i/n$ has a limiting
distribution function $G(x,t)$ and density $g(x,t)$.  If
\[
 F(t):=\int_0^1G(x,t)\,dx,\qquad
 d(t):=\int_0^1g(x,t)\,dx,
\]
assume that $F$ is strictly increasing on the range under consideration and
that $d(t_y)>0$.  With $t_y=F^{-1}(y)$, the natural bulk density is
\begin{equation}\label{eq:general-clock-density}
 \rho_T(x,y)=\frac{g(x,t_y)}{d(t_y)}.
\end{equation}
Thus a bulk local-rank theorem together with
Proposition~\ref{prop:general-clock-tail} gives the same architecture as the
fixed-point proof above, provided one can choose triangular-array parameters
$M_n,B_n,s_n,h_n$ for which the right side of
\eqref{eq:general-clock-tail} tends to zero.  A candidate bulk density is not
enough by itself: one must still prove compensator convergence, or an
equivalent microscopic local law.

Longer cycles require another input.  Their late closing label need not be a
terminal label, so Proposition~\ref{prop:general-clock-tail} alone does not
control them.  In the exponential case, Lemma~\ref{lem:finite-insertion-path}
and Propositions~\ref{prop:cycle-rate-truncation}
and~\ref{prop:cycle-endpoint-bound} overcome this obstruction by swapping a
finite insertion path into the background, charging exceptional rates, and
marking the edge into the largest cycle label.  The finite-path lemma is
distribution-free; the truncation and shell bounds for another clock family
must come from suitable density comparisons and endpoint envelopes.

A concrete testing ground is the common-shape Gamma ranking family of
\cite{stern1990models}.  Shape one is the exponential race and hence the Luce
model.  For fixed Gamma shape and terminal rates bounded below, the terminal
densities have a common polynomial-times-exponential envelope, so
Proposition~\ref{prop:general-clock-tail} applies; a bulk local law remains to
be established.  The independent-utility ranking models of
\cite{thurstone1927law} are another natural target, provided their terminal
densities admit an integrable envelope and the corresponding bulk local limit
is available.  Separately, the finite-dimensional criterion of
\cite{mukherjee2016fixed} already yields the cyclic trace law
\eqref{eq:cycle-intensity} for permuton-sampled permutations with strictly
positive continuous densities, Mallows permutations with $n(1-q_n)$
convergent, and certain continuous exponential-family models.

The extension also has sharp limits.  For each fixed Ewens
parameter $\vartheta>0$, asymptotically uniform one-point marginals coexist
with limiting $\ell$-cycle mean $\vartheta/\ell$~\cite{tavare2021magical};
thus a microscopic hypothesis remains essential.  For fixed $q\in(0,1)$,
Mallows permutations lie in a localized regime in which bounded-cycle counts
have linear means and joint Gaussian rather than Poisson
fluctuations~\cite{he2023cycles}.  These examples underline the central
point: the macroscopic limit predicts the form of the cycle constants only
after a model-specific argument has reached the microscopic scale.

\bibliographystyle{alphaabbr}
\small
\bibliography{fixed_points}

\end{document}
